\documentclass[letterpaper,journal]{IEEEtran}
%\usepackage{ctex}
\usepackage{amsmath,amsfonts}
\usepackage{amsthm}
\usepackage{algorithmic}
\usepackage{algorithm}
\usepackage{array}
\usepackage[caption=false,font=normalsize,labelfont=sf,textfont=sf]{subfig}
\usepackage{textcomp}
\usepackage{stfloats}
\usepackage{url}
\usepackage{verbatim}
\usepackage{graphicx}
\usepackage{epsfig}
% \usepackage{subfigure}
\usepackage{epstopdf}
\usepackage{cite}
\usepackage{booktabs}
\usepackage{multirow} 
\usepackage{amssymb} 
\newtheorem{definition}{Definition}  
\newtheorem{theorem}{Theorem}
\newtheorem{proposition}{Proposition}
\newtheorem*{proposition*}{Proposition}

\newtheorem{assumption}{Assumption}
\newtheorem{lemma}{Lemma}
\newtheorem{corollary}{Corollary}
\newtheorem{remark}{Remark}
%\usepackage{appendix}
\usepackage{placeins}
\usepackage{xr}
\externaldocument{supplemental}
\graphicspath{{./figures/}{./images/}}
\hyphenation{op-tical net-works semi-conduc-tor IEEE-Xplore}

\makeatletter
\newcommand{\rmnum}[1]{\romannumeral #1}
\newcommand{\Rmnum}[1]{\expandafter\@slowromancap\romannumeral #1@}
\makeatother

\renewcommand{\topfraction}{0.99}
\renewcommand{\bottomfraction}{0.99}
\renewcommand{\textfraction}{0.01}
\renewcommand{\floatpagefraction}{0.99}
\renewcommand{\dbltopfraction}{0.9}
\renewcommand{\dblfloatpagefraction}{0.9}
\setlength{\dbltextfloatsep}{8pt plus 2pt}
\setcounter{topnumber}{5}
\setcounter{bottomnumber}{5}
\setcounter{totalnumber}{10}

% Let TeX stretch long unbreakable material (wide equations, long identifiers)
% instead of letting it project a fraction of a point into the margin.
\setlength{\emergencystretch}{2em}

\begin{document}

\title{RAVEN-MCS: Target-Mass-Calibrated Windowed-Asynchronous Federated Data Completion for Mobile Crowdsensing under Two-Stage Non-Random Selection}

\author{**** and Anfeng Liu%
\thanks{This work was supported in part by the National Natural Science Foundation of China (U24A20248).}%
\thanks{*** and Anfeng Liu are with the School of Electronic Information, Central South University, Changsha, 410083, China (E-mail: ****@csu.edu.cn, afengliu@mail.csu.edu.cn).}%
}

%\IEEEpubid{xxxx--xxxx/0x\$xx.00~\copyright~201x IEEE}

\maketitle


\begin{abstract}
Mobile crowdsensing completion is trained on updates selected by
opportunity, observation, and usable-update delivery.  That arrival law
need not match the public deployment target.  RAVEN-MCS treats
windowed-asynchronous federated completion as target-mass alignment.
The server updates once at window close, forms a selection-corrected
reference from legal pre-outcome R/O/U histories, and calibrates it
inside the current feasible window.  When the two laws differ, their
risks are not interchangeable.  The ideal sequential correction recovers
compatible target moments or group mass under support and identification;
it does not prove optimization of deployment-target risk.  Calibration is
uniquely solvable, and the clipped plug-in is a stable perturbation of
that reference when target support is compatible.  The Design factor is an
explicit configuration: RAVEN-Design uses the lagged opportunity/design
ratio, RAVEN-NoDesign leaves that factor at one and does not reconstruct
target mass.  Both settings are reported on every block rather than chosen
by a data-derived rule, because they trade off across endpoints: Design is
worth $1.24\%$ of group mismatch on SensorScope at unchanged prediction, and
on U-Air it costs $6.55\%$ of RMSE while improving cross-sensor mismatch by
$4.99\%$.  A SensorScope development table under
a protocol-constructed Complete-aligned selection reports RAVEN-Design as
structural alignment
with a descriptive $+0.5\%$ worst-group tolerance, not an RMSE win.
On a T-Drive taxi replay,
FedAU-Window and ObsUse-Window remain strongest, and the Design-off setting does
not meet that tolerance versus those two.  No ranking jointly measures
$R$, $O$, and $U$.
\end{abstract}

\begin{IEEEkeywords}
Mobile crowdsensing, federated learning, data completion, windowed asynchronous learning, non-random selection, selection bias, target-mass alignment.
\end{IEEEkeywords}


\section{Introduction}
\label{sec:introduction}

Mobile crowdsensing (MCS) collects urban and environmental data from mobile
devices without a dense dedicated infrastructure
\cite{ganti2011mobile,guo2015mobile,ma2014opportunities}.
Continuous sensing is costly and participation is intermittent, so practical
systems observe only a sparse field and rely on completion
\cite{wang2023outlier,li2024robust,wang2023transformer,wang2024fewshot,wang2025evolutionary,xu2025rethinking,liu2026pyramid,zhang2026mapt}.
Federated learning trains those models without centralizing raw records
\cite{mcmahan2017communication,yang2019federated,kairouz2021advances,li2020federated},
including asynchronous aggregation for irregular or stale updates
\cite{wang2026federated,xie2019asynchronous,liu2023adaptive,sun2024staleness}.
The system studied here is \emph{windowed asynchronous}: updates may arrive at
different times inside a window, but the server applies one global update at
window close.  This is not immediate-update FedAsync.

A less visible issue arises before aggregation: \emph{the data that reach the
server are not necessarily sampled from the distribution that the deployed
predictor is intended to serve}.  Consider an atomic sensing target $i$.  A
mobile client must first have an opportunity to cover $i$; the target must then
actually be observed; finally, the client's local update must complete
computation and communication before the server's usable-update deadline.
Thus, the two-stage MCS process contains three distinct selection layers:
mobile opportunity $R$, observation $O$, and usable update $U$.  Their
probabilities depend on mobility, sensing conditions, resource state, network
quality, deadlines, and staleness.  Even when each individual layer is
operationally reasonable, their composition can make the distribution of
observed and successfully delivered data differ from the public deployment
target.  Optimizing the distribution induced by arriving updates can therefore
optimize the \emph{wrong population risk}.

Existing work leaves an important gap at this point.  Sparse MCS completion
primarily asks how to infer missing values from the available observations
\cite{wang2023spatiotemporal,li2024robust}; asynchronous federated completion
asks how to learn effectively from irregular or stale client updates
\cite{wang2026federated,sun2024staleness}; and participation-aware FL studies
partial, arbitrary, or biased client participation
\cite{wang2022unified,jhunjhunwala2022fedvarp,wu2023anchor,wang2024lightweight,sun2025debiasing,mohri2019agnostic,chen2026fedcure,morishita2026who}.
These directions establish that sparsity, asynchrony, and non-uniform
participation matter, but they do not by themselves resolve the following MCS
problem: \emph{how should one align asynchronous federated data completion with
a deployment target when sensing-side availability and update-side usability
jointly select the data that enter learning?}  In particular, correcting only
the final client-participation distribution does not identify whether the
distortion arose from mobile opportunity, actual observation, or usable-update
censoring; conversely, reconstructing a statistically correct target mass does
not imply that this mass is realizable by the clients available in a finite
mobile window.

We develop \textbf{RAVEN-MCS}, short for
\textbf{Reference Alignment Via Exposure-aware Normalization for Mobile
CrowdSensing}.  RAVEN-MCS treats deployment-target mass alignment, rather than raw
update frequency, as the reference semantics of asynchronous federated data
completion.  As illustrated in Fig.~\ref{fig:raven_overview}, the method first
separates the public deployment target from the arrival/usable-update law.  It
then factors the opportunity, observation, and usable-update layers under
their legal pre-outcome information boundaries (selection-time covariates,
not outcome or test labels).  RAVEN-Design uses the lagged target-mass
opportunity ratio.  When a frozen zero-target heuristic is
positive, RAVEN-NoDesign leaves that factor at one and is not a
target-mass estimator.
The resulting statistical reference may still be unattainable in a particular
mobile window.  RAVEN-MCS therefore maps the corrected reference to a feasible
aggregation decision through instantaneous target calibration, a strongly
convex calibration program, and long-term target debt.
This separation yields three distinct questions:
\emph{what target should be served}, \emph{what target mass the selected data
represent}, and \emph{what target composition can be feasibly realized now}.

\begin{figure*}[t]
\centering
\includegraphics[width=\textwidth]{fig1_raven_problem_framework.pdf}
\caption{RAVEN-MCS aligns asynchronous federated completion with a public
deployment target.  (a)~Three selection layers induce an arrival law that
need not match the deployment target.
(b)~RAVEN-Design forms a selection-corrected target-mass reference and
calibrates it in the current feasible support.  RAVEN-NoDesign uses
$\widehat\zeta=1$ and is not target-mass reconstruction.
(P2) is the per-window calibration program that maps this reference onto the
feasible aggregation set; it is defined in Section~\ref{sec:calibration}.}
\label{fig:raven_overview}
\end{figure*}

Our main contributions are as follows.

\begin{itemize}
    \item \textbf{Deployment-target risk semantics under two-stage non-random
    selection.}
    We formulate asynchronous federated MCS completion as a
    deployment-target alignment problem rather than an arrival-centered
    learning problem.  The formulation distinguishes the atomic deployment law
    from the law induced by mobile opportunity, observation, and usable-update
    selection, and defines structural mismatch at both deployment-group and
    client--opportunity-stratum levels.  We prove that when the deployment and
    arrival laws differ, their induced risk functionals are not generally
    interchangeable.  This provides the formal basis for asking whether a model
    is trained for the population it is actually intended to serve, not whether
    completion RMSE is optimized.

    \item \textbf{Selection-layer-aware target reconstruction.}
    We develop an exposure-aware normalized reconstruction that keeps the
    opportunity/design, observation, and usable-update mechanisms separate
    instead of collapsing them into a single missingness probability.  The
    correction uses only legal pre-outcome information and local
    H\'ajek normalization, and yields a selection-corrected reference mass
    for aggregation.  Usable-update correction is retained as a third
    sequential factor, not as an isolated empirical driver.
    Under compatible target support, sequential pre-outcome ignorability,
    and correct ideal selection ratios, we show that the ideal pooled
    normalized sequential estimator,
    with the target-mass opportunity ratio, recovers compatible target moments
    or target-group mass.  That is not a
    proof that sequential correction optimizes deployment-target risk, nor a
    claim of unbiased local gradients or convergence under asynchrony.
    Clipping is a finite-sample
    stabilizer, not an unbiasedness claim.  That reconstruction is
    RAVEN-Design, the setting in which the Design factor is on.  With the
    factor off, the implemented configuration is RAVEN-NoDesign: the
    opportunity factor is one, and the result is not a target-mass
    reconstruction.  The factor is an explicit configuration rather than a
    data-derived selector, and both settings are reported.  A finite-sample
    coverage bound is audited as an empirical diagnostic in the
    supplemental material.

    \item \textbf{Feasibility-aware instantaneous and long-term target
    calibration.}
    We show that statistically reconstructed target mass and instantaneous
    mobile realizability are different problems.  RAVEN-MCS therefore
    calibrates the corrected reference inside the current feasible aggregation
    set while accounting for target mismatch, reference fidelity, uncertainty,
    and staleness, and carries unresolved group under-service through a
    long-term target-debt process.  The resulting calibration program has a
    nonempty compact convex feasible set and a unique global minimizer under the
    stated conditions.  When the active set is held fixed and target support is
    compatible, the implemented RAVEN-Design clipped plug-in is a Lipschitz perturbation of
    the same-window ideal reference, and the feasibility-aware calibration
    program of Section~\ref{sec:calibration} does not amplify that error.
    Target debt is an internal accounting recursion for residual
    group under-service, not a standalone performance contribution;
    empirically, instantaneous calibration and debt are not prediction
    drivers.
\end{itemize}

RAVEN's novelty is therefore not inverse-probability weighting or H\'ajek
normalization as such.  It is the combination of deployment-target semantics,
selection-layer legal histories, and finite-support calibration.
Isolated usable-update and debt terms are retained for sequential
factorization and accounting; they are not the empirical claim.

We evaluate these contributions with a no-harm control, a risk-semantic
block, and a SensorScope/U-Air development table under always-on protocol
$R$ on station traces (Table~\ref{tab:main_effectiveness}).
Independent hybrid ceilings---a T-Drive taxi replay under real $R$ with
controlled $O$ and $U$, an NSW-Traffic four-region protocol table, a
City Scanner New York {PM}$_{2.5}$ visit-occupancy table
the City Scanner visit-occupancy ranking table in the supplemental material, and a Pilot~2 table with measured
observation the City Scanner Pilot~2 measured-$O$ ranking table in the supplemental material---are not mixed
into that development table.
The Design factor is set explicitly on each block, and both settings are
reported where the trade-off is material
(Table~\ref{tab:main_effectiveness} and Section~\ref{subsec:eval_rq2}).
NSW-Traffic and City Scanner use locally constructed target maps, not a
transferred SensorScope/U-Air lookup object.
On T-Drive, FedAU-Window and
ObsUse-Window remain strongest, and the Design-off setting does not meet the
descriptive $+0.5\%$ worst-group tolerance versus those two.
A supplemental occupancy-robustness block thins protocol $R$ on the same
station clients.  A coverage Gate is audited only as a diagnostic and
does not switch Design.


\section{Related Work}
\label{sec:related_work}

Sparse MCS completion infers unsensed cells from a partial spatiotemporal
field
\cite{ganti2011mobile,guo2015mobile,ma2014opportunities,wang2023spatiotemporal,wang2023outlier,wang2023transformer,wang2024fewshot,li2024robust,wang2025evolutionary,xu2025rethinking,liu2026pyramid,zhang2026mapt,liu2026pricing}.
Recruitment, incentives, and privacy
\cite{liu2026multidimensional,zhang2025incentive,lu2025piece,zhang2026locally}
are orthogonal; they do not formulate deployment-target risk under sequential
selection.

Federated completion keeps sensing records on clients
\cite{mcmahan2017communication,yang2019federated,kairouz2021advances,li2020federated}.
Wang \emph{et al.}~\cite{wang2026federated} combine local completion with a
time-aware architecture and asynchronous aggregation for sparse MCS, in the
spirit of which we study windowed rather than immediate-update asynchrony.
FedAsync and staleness-aware variants update as models arrive
\cite{xie2019asynchronous,liu2023adaptive,sun2024staleness,islam2025seafl,li2026stdfl}.
Those works ask how to learn from the updates that participate.  RAVEN-MCS
additionally models which target mass those updates represent.

Federated conformal prediction builds label-set coverage under
heterogeneity or label shift
\cite{lu2023federated,plassier2023conformal}.  The diagnostic Gate in
Section~\ref{subsec:agg_certificate_alg} is a window-level lower bound on
realized aggregate opportunity coverage $C_t$, not a prediction set and
not the Design rule.

Non-uniform participation is the complementary FL object.  Unified
convergence under arbitrary participation, FedVARP, and anchor sampling
address partial client presence
\cite{wang2022unified,jhunjhunwala2022fedvarp,wu2023anchor}.
FedAU estimates aggregation weights from participation history so that
unknown participation does not replace the original FL objective
\cite{wang2024lightweight}.  Correlated availability and hierarchical
semi-asynchronous bias have likewise been debiased
\cite{sun2025debiasing,chen2026fedcure}; agnostic FL optimizes a worst-case
client mixture \cite{mohri2019agnostic}.  FedProx and SCAFFOLD address
client drift \cite{li2020fedprox,karimireddy2020scaffold}.  These methods
correct participation.  RAVEN-MCS additionally separates sensing-side
opportunity and observation from update-side usability.
Closest in the two-stage FL vocabulary is FedIPW
\cite{morishita2026who}, which factors enrollment then round-level
participation toward a target-population mean client update.  That
factorization is still on the client roster: it does not factor MCS
opportunity or observation of sensing units, and the aligned object is not
a public deployment measure over $h(i)$.  We cite it as the nearest
enrollment--participation IPW, not as an MCS completion baseline.
FLOSS~\cite{goetze2025floss} reweights aggregation for client opt-out and
stragglers.  FLOSS and FedIPW are client-level participation corrections;
the aligned object here is MCS unit-level R/O/U selection of sensing units.
FedPALS~\cite{listozec2024overcoming} aggregates toward a known target
\emph{label} distribution under label shift.  That object is class-label
marginals, not MCS opportunity or observation of sensing units, and the
aligned quantity is not the public deployment measure over $h(i)$; we do
not run it as a comparator.
Classical survey estimators supply the H\'ajek building blocks
\cite{horvitz1952generalization,hajek1971comment,shimodaira2000covariate,robins1994estimation};
importance weighting can lose effect in over-parameterized networks
\cite{byrd2019importance}, which is not claimed here.
A nearest-neighbor comparison on target, selection, sensing-side
correction, and finite-window calibration is in the supplemental material
the nearest-neighbour comparison table in the supplemental material.


\section{System Model and Target-Mass Formulation}
\label{sec:system_model}

\subsection{Windowed Asynchronous Mobile Crowdsensing}
\label{subsec:windowed_mcs}

We consider a mobile crowdsensing (MCS) system with a server and a set of
clients $\mathcal{K}=\{1,\ldots,K\}$.  The common prediction domain is
partitioned into atomic units $i\in\mathcal{I}$; an atomic unit may be written as
$i=(n,t)$, where $n$ is a spatial unit and $t$ is an absolute time index.  Each
unit carries a target value $Y_i\in\mathcal{Y}\subseteq\mathbb{R}$, revealed only
if observed, and a label-free deployment context $C_i^{\mathrm{dep}}$.  All
clients contribute to one completion predictor $f_{\theta}$ with parameters
$\theta\in\Theta\subseteq\mathbb{R}^{d}$, returning
$\widehat Y_i=f_{\theta}(i;C_i^{\mathrm{dep}})\in\mathcal{Y}$ and updated by
local training and aggregation.
A client-side observation of unit $i$ is denoted by $Z_{k,r,i}$ when it is
available to client $k$ in server window $r$.

Server time is divided into aggregation windows
$\mathcal{W}_r=[T_r,T_{r+1})$ for $r=0,1,\ldots$.
The server keeps the parameter vector $\theta_r\in\Theta$ fixed during
$\mathcal{W}_r$.  Client $k$ may start local training from a previously
downloaded model $\theta_{\sigma_{k,r}}$,
where $\sigma_{k,r}\le r$.  We use
$\tau_{k,r}\triangleq r-\sigma_{k,r}$
as the model-version staleness of that attempt.  Updates may arrive at
different times within one window, but individual arrivals do not
immediately modify the global model.  The server closes the window,
determines which registered attempts are usable, computes one
aggregation vector, and then performs the single transition
$\theta_r\rightarrow\theta_{r+1}$.  This \emph{windowed-asynchronous}
semantics is distinct from immediate-update FedAsync, in which an
arrival would revise the global model before the window closes.  It
keeps exogenous mobile/update events independent of the method-specific
aggregation weights.

\subsection{Two-Stage, Three-Layer Non-Random Selection}
\label{subsec:selection_process}

The title's \emph{two-stage selection} consists of a sensing-side stage and an
update-side stage.  The sensing-side stage itself has two layers: mobile
opportunity and actual observation.  The update-side stage determines whether a
registered client attempt becomes usable before window closure.

For client $k$, window $r$, and unit $i$, let
$R_{k,r,i}\in\{0,1\}$ be the realized mobile opportunity indicator.
It indicates whether $(k,i)$ belongs to the opportunity/risk set in
window $r$ before the current sensing outcome is revealed.
Conditional on $R_{k,r,i}=1$,
$O_{k,r,i}\in\{0,1\}$
indicates whether unit $i$ is actually observed.  A client with a
nonempty observed buffer registers an update attempt; the set of
registered clients is $\mathcal{E}_r$.  For $k\in\mathcal{E}_r$, the
attempted update is usable ($U_{k,r}=1$) only if computation and
transmission succeed, the update arrives by $T_{r+1}$, and its
staleness satisfies $\tau_{k,r}\le S_{\max}$, where $S_{\max}$ is the block's
frozen maximum admissible model-version staleness.  Otherwise,
$U_{k,r}=0$.  We first define the usable set
$\mathcal{U}_r\triangleq\{k\in\mathcal{E}_r:U_{k,r}=1\}$.
After selection-corrected masses are formed in
Section~\ref{sec:reconstruction}, zero-mass clients are removed and the
actual aggregation set $\mathcal{A}_r\subseteq\mathcal{U}_r$ is formed.
Failures, timeouts, network interruptions, and stale arrivals remain
recorded as $U_{k,r}=0$ rather than disappearing from the selection
log.

We distinguish two structural mappings from units to finite label-free index
sets: the \emph{deployment-target group} $h:\mathcal{I}\to\mathcal{G}$ and the
\emph{opportunity stratum} $s:\mathcal{I}\to\mathcal{S}$, with $\mathcal{G}$ the
finite set of deployment groups and $\mathcal{S}$ the finite set of
opportunity strata.  The group mapping specifies which deployment
mass the predictor should serve, whereas the stratum mapping pools
repeatable mobile exposure patterns for selection modeling.  The two
mappings are different scientific objects, so $h(i)\neq s(i)$ in
general.  Both mappings are frozen from public/training-side
information and do not use test-label values.

For theoretical bookkeeping, let the normalized opportunity-event law in
window $r$ be
\begin{equation}
 P_r^{\mathrm{opp}}(k,i)
 =
 \pi^{\mathrm{opp}}_{k,s(i),r}
 \nu^{\mathrm{opp}}_{k,r}(i\mid s(i)),
 \qquad
 \sum_{k,i}P_r^{\mathrm{opp}}(k,i)=1,
 \label{eq:opp_event_law}
\end{equation}
where $\pi^{\mathrm{opp}}_{k,s,r}$ is the client--stratum opportunity mass and
$\nu^{\mathrm{opp}}_{k,r}(i\mid s)$ is the conditional atomic opportunity law.
The realized risk-set indicators $R_{k,r,i}$ are the implementation-level
bookkeeping of the opportunity events induced by this mobile exposure process;
they are not assumed to be mutually independent Bernoulli draws.

The subsequent selection probabilities are conditional probabilities under the
same selection process, with $\Pr\{\cdot\}$ that probability.
\begin{align}
 p^{\mathrm{obs}}_{k,r,i}
 &\triangleq
 \Pr\!\left(
 O_{k,r,i}=1
 \mid
 R_{k,r,i}=1,\mathcal{H}^{\mathrm{obs}}_{k,r,i}
 \right),
 \label{eq:pobs_def}\\
 q^{\mathrm{use}}_{k,r}
 &\triangleq
 \Pr\!\left(
 U_{k,r}=1
 \mid
 k\in\mathcal{E}_r,\mathcal{H}^{\mathrm{use}}_{k,r}
 \right).
 \label{eq:quse_def}
\end{align}
The observation history $\mathcal{H}^{\mathrm{obs}}_{k,r,i}$ contains
only information available before $O_{k,r,i}$ is realized.  The
usable-update history $\mathcal{H}^{\mathrm{use}}_{k,r}$ is measured
after the sensing-side stage but before $U_{k,r}$ is realized.  It may
therefore contain the already observed registration state, but not the
realized computation/transmission outcomes, future events, or the
current update vector used for aggregation.  This sequential
information boundary is part of the statistical model.

\subsection{Public Deployment Target and Compatible Target Measures}
\label{subsec:target_measures}

We begin with a public, label-independent atomic deployment target over the
deployment domain $\mathcal{I}_{\mathrm{dep}}\subseteq\mathcal{I}$,
\begin{equation}
    \varpi_i^{\mathrm{tar}}\ge0,
    \qquad
    \sum_{i\in\mathcal{I}_{\mathrm{dep}}}
    \varpi_i^{\mathrm{tar}}=1.
    \label{eq:atomic_target}
\end{equation}
Its deployment-group marginal is
$\mu_g=\sum_{i:h(i)=g}\varpi_i^{\mathrm{tar}}$, so
$\boldsymbol{\mu}\in\Delta^{|\mathcal{G}|-1}$.
These quantities define the application-level target served by the deployed
predictor.  Opportunity correction uses a compatible client--stratum target
$\pi^{\mathrm{tar}}_{k,i}$ built from the stratum mass of $s(i)$, a frozen
public client share, and the conditional atomic law on $s(i)$
(supplemental material).  Deployment groups $h(i)$ and opportunity strata
$s(i)$ need not coincide.  Compatibility requires
\begin{equation}
    \sum_k\pi^{\mathrm{tar}}_{k,i}
    =
    \varpi_i^{\mathrm{tar}},
    \qquad
    \sum_{k,s}\pi^{\mathrm{tar}}_{k,s}=1.
    \label{eq:target_compatibility}
\end{equation}
Equation~\eqref{eq:target_compatibility} is the bridge between the
deployment-risk target and the client--stratum target used for Design
and coverage.  Define $I_{\mathrm{tar}}=\{(k,s):\pi^{\mathrm{tar}}_{k,s}>0\}$.
The unit lookup used by the target-coverage statistic of the supplemental
material, and the pair opportunity mass used by
$C_t$, are constructed in the supplemental material.
All target quantities are frozen before test-label evaluation.

\subsection{Aggregate Coverage Certification Object}
\label{subsec:agg_coverage_object}

Realized aggregate coverage is
\begin{equation}
 C_t
 =
 \sum_{(k,s)\in I_{\mathrm{tar}}}
 \pi^{\mathrm{tar}}_{k,s}
 \,
 \mathbf{1}\!\left\{\pi^{\mathrm{opp}}_{k,s,t}\ge\pi_{\mathrm{safe}}\right\},
 \label{eq:agg_Ct}
\end{equation}
where $\pi^{\mathrm{opp}}_{k,s,t}$ is the lagged opportunity mass of
that pair, $\pi_{\mathrm{safe}}$ is a frozen threshold, and
$\mathbf{1}\{\cdot\}$ is the indicator of the event in braces.
A diagnostic Gate on a lower bound of $C_t$ does not switch Design.

\subsection{Deployment-Target Risk Versus Arrival Risk}
\label{subsec:target_risk}

For the squared loss
$\ell(\widehat Y,Y)=\tfrac{1}{2}(\widehat Y-Y)^2$, the deployment-target risk is
\begin{equation}
 \mathcal{R}_{\mu}(\theta)
 =
 \sum_{i\in\mathcal{I}_{\mathrm{dep}}}
 \varpi_i^{\mathrm{tar}}
 \ell\!\left(
 f_{\theta}(i;C_i^{\mathrm{dep}}),Y_i
 \right).
 \label{eq:target_risk}
\end{equation}
By contrast, the data reaching learning are distributed according to the
opportunity, observation, and usable-update mechanisms.

For evaluation, quantities without a window index denote the frozen
evaluation-period marginals of their corresponding pre-outcome mechanisms.
Under \eqref{eq:opp_event_law}--\eqref{eq:quse_def}, the joint sequential
probability that client $k$ contributes unit $i$ to learning is
$P^{\mathrm{opp}}(k,i)\,p^{\mathrm{obs}}_{k,i}\,q^{\mathrm{use}}_{k}$.
The label-free atomic arrival intensity is the corresponding sum, and the
evaluation plug-in is
\begin{equation}
 \widehat r_i^{\mathrm{arr}}
 =
 \sum_k
 \widehat\pi^{\mathrm{opp}}_{k,s(i)}
 \widehat\nu^{\mathrm{opp}}_{k}(i\mid s(i))
 \widehat p^{\mathrm{obs}}_{k,i}
 \widehat q^{\mathrm{use}}_{k,s(i)},
 \label{eq:arrival_intensity}
\end{equation}
where $\widehat q^{\mathrm{use}}_{k,s}$ is the corresponding
evaluation-period marginal of the legal pre-$U$ usable-update model over
registered attempts associated with stratum $s$.  Normalize
\begin{equation}
 \widetilde\rho_i^{\mathrm{arr}}
 =
 \frac{\widehat r_i^{\mathrm{arr}}}
 {\sum_{j\in\mathcal{I}_{\mathrm{dep}}}\widehat r_j^{\mathrm{arr}}},
 \label{eq:arrival_distribution}
\end{equation}
and define
\begin{equation}
 \mathcal{R}_{\rho}(\theta)
 =
 \sum_{i\in\mathcal{I}_{\mathrm{dep}}}
 \widetilde\rho_i^{\mathrm{arr}}
 \ell\!\left(
 f_{\theta}(i;C_i^{\mathrm{dep}}),Y_i
 \right).
 \label{eq:arrival_risk}
\end{equation}
The two risks answer different questions.  The deployment-target risk
$\mathcal{R}_{\mu}$ measures how accurately the predictor serves the deployment
population it is intended for; the arrival risk $\mathcal{R}_{\rho}$ measures how
accurately it fits the updates that actually reach training.  The loss and the
predictor are the same in both; only the weighting law differs.  Since
opportunity, observation, and usable-update selection are non-random, those two
populations need not coincide, so accuracy on the arriving updates need not
transfer to the deployment population.  RAVEN-MCS therefore treats
deployment-target mass alignment, rather than usable-update frequency, as the
reference semantics of the learning problem.

\subsection{Structural Alignment Objective}
\label{subsec:design_objective}

Section~\ref{sec:reconstruction} will construct two corrected within-client
compositions: a deployment-group composition
$\mathbf{c}^{g}_{k,r}\in\Delta^{|\mathcal{G}|-1}$ and an opportunity-stratum
composition $\mathbf{c}^{s}_{k,r}\in\Delta^{|\mathcal{S}|-1}$.  For the final
aggregation set $\mathcal{A}_r$, define
$M_r=[\mathbf{c}^{g}_{k,r}]_{k\in\mathcal{A}_r}$.
Given feasible aggregation weights $\boldsymbol\alpha_r\in\mathcal{C}_r$, the
effective deployment-group composition in an active window is
$\boldsymbol\omega_r=M_r\boldsymbol\alpha_r$.  The feasible set
$\mathcal{C}_r$ is \eqref{eq:Cr}, and $\boldsymbol\alpha_r$ is determined by the
calibration program (P2) of Section~\ref{subsec:P2} rather than supplied as an
input.

Let $\eta_r>0$ denote the server step size and define
$I_r=\mathbf{1}\{|\mathcal{A}_r|>0\}$ and
$S_R=\sum_{r=0}^{R-1}\eta_r I_r$.
For $S_R>0$,
\begin{equation}
    \overline{\boldsymbol\omega}_R
    =
    \frac{1}{S_R}
    \sum_{r=0}^{R-1}
    \eta_r I_r\boldsymbol\omega_r.
    \label{eq:omega_bar}
\end{equation}
We measure the group-level structural mismatch
$\Delta_{\mathrm{group}}(R)=\left\|
\overline{\boldsymbol\omega}_R-\boldsymbol\mu
\right\|_1$,
the $\ell_1$ distance between the time-averaged effective deployment-group
composition and the deployment target.  Both sides are probability vectors
over $\mathcal{G}$, so $0\le\Delta_{\mathrm{group}}\le2$.

For client--stratum structure, extend $\alpha_{k,r}=0$ to
$k\notin\mathcal{A}_r$ and define
\begin{equation}
 \widehat\Pi_{k,s}(R)
 =
 \frac{1}{S_R}
 \sum_{r=0}^{R-1}
 \eta_r I_r\,
 \alpha_{k,r}\,
 c^{s}_{k,r,s}.
 \label{eq:Pi_hat_ks}
\end{equation}
Because each active-window $\boldsymbol\alpha_r$ and each
$\mathbf{c}^{s}_{k,r}$ are probability vectors,
$\sum_{k,s}\widehat\Pi_{k,s}(R)=1$.  We therefore measure
\begin{equation}
 \Delta_{c\text{-}s}(R)
 =
 \sum_{k,s}
 \left|
 \widehat\Pi_{k,s}(R)-\pi^{\mathrm{tar}}_{k,s}
 \right|.
 \label{eq:delta_cs}
\end{equation}
The overall design goal is to keep $\mathcal{R}_{\mu}$ low while
reducing $\Delta_{\mathrm{group}}$ and $\Delta_{c\text{-}s}$ under
time-varying mobile feasibility and legal pre-outcome information.


\section{Selection-Layer-Aware Target Reconstruction}
\label{sec:reconstruction}

RAVEN-MCS (Reference Alignment Via Exposure-aware Normalization) first forms a
selection-corrected reference for the mass the usable data represent.
Opportunity, observation, and usable-update selection are kept separate
because they have different conditioning sets and therefore cannot be
replaced by a single generic missingness probability.

\subsection{Opportunity/Design Correction}
\label{subsec:design_correction}

For a target-supported client--atomic pair, the ideal opportunity/design ratio is
\begin{equation}
 \zeta^{o}_{k,r,i}
 =
 \frac{\pi^{\mathrm{tar}}_{k,s(i)}}
      {\pi^{\mathrm{opp}}_{k,s(i),r}}
 \frac{\nu^{\mathrm{tar}}_{i\mid s(i)}}
      {\nu^{\mathrm{opp}}_{k,r}(i\mid s(i))}
 =
 \frac{\pi^{\mathrm{tar}}_{k,i}}
      {P_r^{\mathrm{opp}}(k,i)}.
 \label{eq:design_ratio_general}
\end{equation}
Under within-stratum exchangeability,
$\nu^{\mathrm{opp}}_{k,r}(i\mid s)=\nu^{\mathrm{tar}}_{i\mid s}$, this reduces to
\begin{equation}
 \widehat\zeta^{o}_{k,r,s}
 =
 \frac{\pi^{\mathrm{tar}}_{k,s}}
 {\max\{\widehat\pi^{\mathrm{opp}}_{k,s,r},\pi_{\min}\}}.
 \label{eq:design_ratio_simple}
\end{equation}
The opportunity estimate is lagged: current-window outcomes cannot alter the
weight applied to that same window.  One admissible implementation uses
\begin{align*}
 C^{\mathrm{opp}}_{k,s,r}
 &=
 \rho_{\mathrm{opp}}C^{\mathrm{opp}}_{k,s,r-1}
 +N^{\mathrm{opp}}_{k,s,r-1},\\
 \widehat\pi^{\mathrm{opp}}_{k,s,r}
 &=
 \frac{C^{\mathrm{opp}}_{k,s,r}+\epsilon_{\mathrm{opp}}}
 {\sum_{j,u}
 (C^{\mathrm{opp}}_{j,u,r}+\epsilon_{\mathrm{opp}})}.
\end{align*}
Other pre-outcome estimators are compatible with the framework if they preserve
the same support and information boundary.

The Design factor $D_t\in\{\mathrm{on},\mathrm{off}\}$ of
\eqref{eq:design_factor} applies \eqref{eq:design_ratio_simple} when it is
on, and that setting is RAVEN-Design.
With the factor off,
Algorithm~\ref{alg:window_close} runs RAVEN-NoDesign and sets $\widehat\zeta^{o}_{k,r,i}=1$ on
observed pairs.  The factor is a configuration of the deployment, supplied to
Algorithm~\ref{alg:window_close} before the window opens, not a statistic
computed from the data; the same two settings are reported on every block.
The off setting is an uncorrected opportunity factor: it
does not replace $\pi^{\mathrm{opp}}$ by $\pi^{\mathrm{tar}}$, and it does
not reconstruct target mass on the opportunity/design layer.

\subsection{Observation Correction and Local H\'ajek Normalization}
\label{subsec:observation_hajek}

The ideal first-stage weight for a realized opportunity is
$a^\star_{k,r,i}=\zeta^{o}_{k,r,i}/p^{\mathrm{obs}}_{k,r,i}$.
The implemented plug-in weight is stabilized as
\begin{equation}
 a_{k,r,i}
 =
 \min\!\left\{
 a_{\max},
 \frac{\widehat\zeta^{o}_{k,r,i}}
 {\max\{\widehat p^{\mathrm{obs}}_{k,r,i},p_{\min}\}}
 \right\}.
 \label{eq:a_stabilized}
\end{equation}
Define the corrected total mass
\begin{equation}
    m_{k,r}
    =
    \sum_i
    R_{k,r,i}O_{k,r,i}a_{k,r,i},
    \label{eq:m_total}
\end{equation}
the corrected deployment-group mass
\begin{equation}
 m^{g}_{k,r,g}
 =
 \sum_i
 R_{k,r,i}O_{k,r,i}a_{k,r,i}
 \mathbf{1}\{h(i)=g\},
 \label{eq:m_group}
\end{equation}
and the corrected opportunity-stratum mass
\begin{equation}
 m^{s}_{k,r,s}
 =
 \sum_i
 R_{k,r,i}O_{k,r,i}a_{k,r,i}
 \mathbf{1}\{s(i)=s\}.
 \label{eq:m_stratum}
\end{equation}
RAVEN aggregates only target-relevant usable clients:
$\mathcal{A}_r\triangleq\{k\in\mathcal{U}_r:m_{k,r}>0\}$.
Every $k\in\mathcal{A}_r$ has well-defined normalized record weights
$\overline a_{k,r,i}=\frac{R_{k,r,i}O_{k,r,i}a_{k,r,i}}{m_{k,r}}$.
Its deployment-group composition is
$c^{g}_{k,r,g}=\frac{m^{g}_{k,r,g}}{m_{k,r}}$, and its
opportunity-stratum composition is
$c^{s}_{k,r,s}=\frac{m^{s}_{k,r,s}}{m_{k,r}}$.
Both composition vectors sum to one for $k\in\mathcal{A}_r$.

The local effective sample size is
\begin{equation}
 n^{\mathrm{eff}}_{k,r}
 =
 \frac{m_{k,r}^2}
 {\sum_i
  (R_{k,r,i}O_{k,r,i}a_{k,r,i})^2}
 =
 \frac{1}{\sum_i\overline a_{k,r,i}^2}.
 \label{eq:neff}
\end{equation}
It measures weight concentration and is deliberately distinct from the
target-equivalent corrected mass $m_{k,r}$.

The local completion objective is
\begin{equation}
 \widehat F^{H}_{k,r}(\theta)
 =
 \sum_i
 \overline a_{k,r,i}
 \frac{1}{2}
 \left(
 f_{\theta}(i;C_i^{\mathrm{dep}})
 -
 Z_{k,r,i}
 \right)^2.
 \label{eq:local_hajek_objective}
\end{equation}
Starting from $\theta_{\sigma_{k,r}}$, client $k$ performs
$E_{\mathrm{loc}}$ local steps of $\ell$ with step size $\gamma_r>0$, reaching
$\theta^{(E_{\mathrm{loc}})}_{k,r}$.  Its normalized model
displacement is
\begin{equation}
 \mathbf{u}_{k,r}
 =
 \frac{
 \theta_{\sigma_{k,r}}
 -
 \theta^{(E_{\mathrm{loc}})}_{k,r}
 }{\gamma_r E_{\mathrm{loc}}}.
 \label{eq:normalized_update}
\end{equation}

\subsection{Usable-Update Correction and Reference Aggregation Mass}
\label{subsec:usable_correction}

For a registered attempt, $\widehat q^{\mathrm{use}}_{k,r}$ is estimated only
from variables measurable before $U_{k,r}$ is realized.  For
$k\in\mathcal{A}_r$, define
\begin{equation}
 d_{k,r}
 =
 \min\!\left\{
 d_{\max},
 \frac{1}
 {\max\{\widehat q^{\mathrm{use}}_{k,r},q_{\min}\}}
 \right\},
 \label{eq:d_use}
\end{equation}
and $b_{k,r}=m_{k,r}d_{k,r}$.
The selection-corrected reference aggregation vector is
\begin{equation}
 \widehat\beta_{k,r}
 =
 \frac{b_{k,r}}
 {\sum_{j\in\mathcal{A}_r}b_{j,r}},
 \qquad
 k\in\mathcal{A}_r.
 \label{eq:beta_ref}
\end{equation}
Thus $\widehat{\boldsymbol\beta}_r$ is a probability vector on
$\mathcal{A}_r$.  The ideal untruncated factorization of
$\beta^\star\overline a^\star$ is in the supplemental material; the
implemented $\widehat{\boldsymbol\beta}_r$ is a stabilized plug-in and is
not claimed to preserve exact unbiasedness when clipping is active.
The corrected deployment-group composition represented by the usable
data before server calibration is
$\boldsymbol\omega_r^{\mathrm{ref}}=M_r\widehat{\boldsymbol\beta}_r$.
Even a statistically valid reconstruction need not satisfy
$\boldsymbol\omega_r^{\mathrm{ref}}=\boldsymbol\mu$ in a finite mobile
window.  The current feasible client set can lack target support.

\subsection{Information Boundary and Stabilization}
\label{subsec:information_boundary}

Every correction applied in window $r$ is measurable before its own random
outcome.  Opportunity estimates use lagged opportunity history,
observation propensities exclude the current sensing outcome and label,
and usable-update propensities exclude the realized $U$ outcome and
future events.  The server determines $\boldsymbol\alpha_r$ without
inspecting the coordinates, norm, or direction of a current random
update for the purpose of assigning that update's weight.  The clipping
constants in \eqref{eq:a_stabilized} and \eqref{eq:d_use} are
finite-sample bias--variance stabilizers.  Exact identification in
Section~\ref{sec:theory} concerns the ideal ratios, or equivalently
regimes in which the stabilizing bounds are inactive on target support.

\subsection{Why Pairwise Positive Lower Bounds Are Not the Formal Certificate}
\label{subsec:pairwise_barrier}

A deterministic pre-window pairwise lower bound is structurally vacuous
under hidden zero-opportunity events by
the positive pairwise lower-bound barrier proposition in the supplemental material;
the availability-induced barrier corollary of the supplemental material places the frozen stress levels
$a\in\{0.80,0.60,0.40,0.20\}$ in that region.  Pairwise positive LCBs are
not the aggregate object \eqref{eq:agg_Ct} and do not switch Design.


\section{Feasibility-Aware Target Calibration}
\label{sec:calibration}

Selection reconstruction determines what the usable data statistically
represent.  It does not guarantee that the intended deployment target is
instantaneously realizable from the currently available clients.  RAVEN
therefore treats $\widehat{\boldsymbol\beta}_r$ as a corrected reference and
calibrates it within the support feasible at window closure.

\subsection{Feasible Aggregation Set}
\label{subsec:feasible_set}

Let $K_r=|\mathcal{A}_r|$ in an active window and retain
$M_r=[\mathbf{c}^{g}_{k,r}]_{k\in\mathcal{A}_r}$.  The server chooses
$\boldsymbol\alpha_r\in\mathbb{R}^{K_r}$ from
\begin{equation}
 \mathcal{C}_r
 =
 \left\{
 \boldsymbol\alpha:
 \mathbf{1}^{\top}\boldsymbol\alpha=1,\;
 0\le\alpha_k\le\overline\alpha_r,\;
 \|\boldsymbol\alpha\|_2^2\le\frac{1}{\overline E_r}
 \right\},
 \label{eq:Cr}
\end{equation}
where the coordinate cap $\overline\alpha_r$ and $\ell_2$ support bound
$\overline E_r$ are frozen functions of $K_r$
the effective-support cap of the supplemental material.
Here $0<\alpha_{\max}\le1$ and $E_{\min}\ge1$ are frozen design parameters.
The coordinate cap controls single-client dominance, while the
$\ell_2$ constraint imposes a minimum effective aggregation support.  The
uniform vector $\mathbf{1}/K_r$ belongs to $\mathcal{C}_r$, so every active
window has at least one feasible aggregation vector.

RAVEN also uses pre-aggregation uncertainty only.  A lagged variance proxy is
\begin{equation}
 v_{k,r}
 =
 \frac{
 \overline S_{k,r^-}^{\,2}
 }{
 \max\{n^{\mathrm{eff}}_{k,r},1\}
 }
 +
 v_{\mathrm{floor}},
 \qquad
 V_r=\operatorname{diag}(v_{k,r})\succeq0,
 \label{eq:variance_proxy}
\end{equation}
where $\overline S_{k,r^-}^{\,2}$ is available before the current random update
is weighted.  Let $\overline{\boldsymbol\tau}_r\ge0$ denote the frozen
pre-aggregation staleness vector.

\subsection{Long-Term Target Debt}
\label{subsec:debt}

Finite-window feasibility can make
$M_r\boldsymbol\alpha=\boldsymbol\mu$ impossible.  RAVEN maintains a
nonnegative target-debt vector
$\mathbf{Q}_r\in\mathbb{R}^{|\mathcal{G}|}_+$, initialized by
$\mathbf{Q}_0=\mathbf{0}$.  In an active window,
$\mathbf{Q}_{r+1}=\left[\mathbf{Q}_r+\eta_r(\boldsymbol\mu-\boldsymbol\omega_r)\right]_+$,
with $\boldsymbol\omega_r=M_r\boldsymbol\alpha_r$.
If the window is inactive, $\mathbf{Q}_{r+1}=\mathbf{Q}_r$.  A positive
$Q_{r,g}$ records accumulated under-service of deployment group $g$ without
pretending that unavailable target mass can be synthesized in the current
window.

For every prefix with $S_R>0$, the recursion implies the deterministic
accounting relation
\begin{equation}
 \left\|
 \overline{\boldsymbol\omega}_R-\boldsymbol\mu
 \right\|_1
 \le
 \frac{
 2\|\mathbf{Q}_R\|_1
 }{S_R}.
 \label{eq:debt_bound}
\end{equation}
The full derivation is in the target-debt accounting proof of the supplemental material.

\subsection{Feasibility-Aware Calibration Program}
\label{subsec:P2}

For every active window, RAVEN solves
$\text{(P2)}\ \min_{\boldsymbol\alpha\in\mathcal{C}_r}J_r(\boldsymbol\alpha)$.
The objective is
\begin{align}
 J_r(\boldsymbol\alpha)
 ={}&-\mathbf{Q}_r^{\top}M_r\boldsymbol\alpha
 +\frac{\lambda_g}{2}
 \|M_r\boldsymbol\alpha-\boldsymbol\mu\|_2^2
 \nonumber\\
 &+\frac{\lambda_{\beta}}{2}
 \|\boldsymbol\alpha-\widehat{\boldsymbol\beta}_r\|_2^2
 \nonumber\\
 &+\frac{\lambda_v}{2}
 \boldsymbol\alpha^{\top}V_r\boldsymbol\alpha
 +\lambda_s\overline{\boldsymbol\tau}_r^{\top}\boldsymbol\alpha.
\label{eq:P2_objective}
\end{align}
The debt term prioritizes accumulated under-served groups; the group,
reference, variance, and staleness terms are coordinated terms of one
program.  $\widehat{\boldsymbol\beta}_r$ is the selection-corrected
reference; $\boldsymbol\alpha_r$ is the feasible aggregation decision.

\subsection{Global Update and Window Logic}
\label{subsec:global_update}

Let $\boldsymbol\alpha_r^\star$ denote the solution of (P2).  The server update
is
\begin{equation}
 \theta_{r+1}
 =
 \theta_r
 -
 \eta_r
 \sum_{k\in\mathcal{A}_r}
 \alpha^\star_{k,r}\mathbf{u}_{k,r}.
 \label{eq:server_update}
\end{equation}
After aggregation, the server updates the debt via
$\mathbf{Q}_{r+1}=\left[\mathbf{Q}_r+\eta_r(\boldsymbol\mu-\boldsymbol\omega_r)\right]_+$,
where $\boldsymbol\omega_r=M_r\boldsymbol\alpha_r$, and advances lagged
selection and variance states.  If $\mathcal{A}_r=\varnothing$, the
window is inactive: $\theta_{r+1}=\theta_r$, $\mathbf{Q}_{r+1}=\mathbf{Q}_r$,
and the window does not enter $S_R$.  Algorithm~\ref{alg:window_close}
records this order: reconstruct $\widehat{\boldsymbol\beta}_r$, solve (P2) for
$\boldsymbol\alpha_r^\star$, then aggregate and update debt.

\subsection{Finite-Sample Aggregate Coverage Certificate}
\label{subsec:agg_certificate_alg}

The Design factor $\eqref{eq:design_factor}$ is configured on or off.  On, it
is RAVEN-Design and uses lagged opportunity ratios; off, it is RAVEN-NoDesign
and uses the uncorrected factor $1$.  A diagnostic Gate,
detailed in the supplemental material as the aggregate coverage certificate algorithm of the supplemental material,
uses only pre-window information $F_{t-}$ with
\begin{equation}
 k_n=\lceil(n+1)(1-\alpha)\rceil,
 \qquad
 \alpha=0.05,
 \label{eq:agg_kn}
\end{equation}
and
\begin{equation}
 C^{\mathrm{LCB}}_t=\max\{0,\widehat C_t-q_t\},
 \label{eq:agg_CLCB}
\end{equation}
so $\mathrm{Gate}_t=\mathbf{1}\{C^{\mathrm{LCB}}_t\ge\tau_{\mathrm{cov}}\}$ with
$\tau_{\mathrm{cov}}=0.05$.  WARMUP and RESET force the Gate off.  It does
not switch Design.

Let $D_t\in\{\mathrm{on},\mathrm{off}\}$ be the Design factor of window $t$:
a configuration of the deployment, fixed before the window opens and recorded
with it.  The implemented opportunity/design factor is
\begin{equation}
\widehat\zeta^{o}_{k,t,i}
=
\begin{cases}
\dfrac{\pi^{\mathrm{tar}}_{k,s(i)}}
{\max\{\widehat\pi^{\mathrm{opp}}_{k,t,s(i)},\pi_{\min}\}}, & D_t=\mathrm{on},\\[1.5mm]
1, & D_t=\mathrm{off}.
\end{cases}
\label{eq:design_factor}
\end{equation}
With $D_t=\mathrm{on}$ this is the lagged ratio
\eqref{eq:design_ratio_simple}, the setting called RAVEN-Design; with
$D_t=\mathrm{off}$ it is the uncorrected factor one, called RAVEN-NoDesign,
and the result is not target-mass reconstruction.
\eqref{eq:design_factor} is the only place the two settings differ: the
observation, usable-update, calibration, and debt steps are identical.
We report both settings on every block instead of selecting one from the
data, because the choice trades off across endpoints
(Section~\ref{subsec:eval_rq2}), and a data-derived selector would have to be
justified on the same evaluation it changes.
As a descriptive statistic, and not as a selector, the supplemental material
reports the share of training risk units whose attached client--stratum
carries zero public target mass.

\begin{algorithm}[t]
\caption{Window-close RAVEN-MCS with the Design factor as an explicit configuration}
\label{alg:window_close}
\begin{algorithmic}[1]
\REQUIRE window $t$; recorded $R_{k,t,i}$, $O_{k,t,i}$, $U_{k,t}$;
public $\boldsymbol\mu$; debt $\mathbf{Q}_t$; model $\theta_t$;
frozen public $\pi^{\mathrm{tar}}$; Design factor $D_t\in\{\mathrm{on},\mathrm{off}\}$
\ENSURE $\theta_{t+1}$, $\mathbf{Q}_{t+1}$
\STATE Form $\mathcal{E}_t$ and
$\mathcal{U}_t=\{k\in\mathcal{E}_t:U_{k,t}=1\}$.
\IF{$D_t=\mathrm{on}$}
\STATE Compute lagged opportunity/design ratios
$\widehat\zeta^{o}_{k,t,i}$ \COMMENT{RAVEN-Design}
\ELSE
\STATE Set $\widehat\zeta^{o}_{k,t,i}\leftarrow 1$ on observed pairs \COMMENT{RAVEN-NoDesign}
\ENDIF
\STATE Compute observation weights $a_{k,t,i}$ and masses $m_{k,t}$;
set $\mathcal{A}_t=\{k\in\mathcal{U}_t:m_{k,t}>0\}$.
\IF{$\mathcal{A}_t=\varnothing$}
\STATE Set $\theta_{t+1}\leftarrow\theta_t$,
$\mathbf{Q}_{t+1}\leftarrow\mathbf{Q}_t$, and stop.
\ELSE
\STATE Compute usable-update factors $d_{k,t}$ and reference
$\widehat{\boldsymbol\beta}_t$.
\STATE Form $\mathcal{C}_t$ and lagged $V_t$,
$\overline{\boldsymbol\tau}_t$; solve (P2) for
$\boldsymbol\alpha_t^\star$.
\STATE $\theta_{t+1}\leftarrow\theta_t
-\eta_t\sum_{k\in\mathcal{A}_t}\alpha^\star_{k,t}\mathbf{u}_{k,t}$.
\STATE $\boldsymbol\omega_t\leftarrow M_t\boldsymbol\alpha_t^\star$,
$\mathbf{Q}_{t+1}\leftarrow
[\mathbf{Q}_t+\eta_t(\boldsymbol\mu-\boldsymbol\omega_t)]_+$.
\STATE Advance lagged selection and variance states.
\ENDIF
\end{algorithmic}
\end{algorithm}


\section{Theoretical Analysis}
\label{sec:theory}

We state the results used by the mainline.  Proposition~\ref{prop:reconstruction}
recovers compatible target moments or group mass under A1--A3 for the
ideal Design-on sequential estimator; it does not
prove optimization of deployment-target risk, unbiased local gradients, or
convergence under asynchrony.

\subsection{Assumptions}
\label{subsec:assumptions}

\noindent\textbf{A1 (Compatible target support and positivity).}
Public targets satisfy \eqref{eq:atomic_target} and
\eqref{eq:target_compatibility}.  Inverse factors are applied only where
$P_r^{\mathrm{opp}}$, $p^{\mathrm{obs}}$, and $q^{\mathrm{use}}$ are
strictly positive on target-supported events.

\noindent\textbf{A2 (Sequential pre-outcome ignorability).}
Conditional on the legal history at each stage, $O$ and $U$ are
conditionally mean-ignorable for the reconstructed moment.  The $O$ model
uses only $\mathcal{H}^{\mathrm{obs}}$; the $U$ model may use the realized
sensing-side registration state, but not the future $U$ outcome
(display in the supplemental material).

\noindent\textbf{A3 (Correct ideal ratios and cluster-level regularity).}
Ideal ratios equal their population counterparts, or stabilizing bounds are
inactive on relevant support.  As $N_R^{\mathrm{opp}}\rightarrow\infty$,
weighted cluster sums obey a weak law and the H\'ajek denominator stays
strictly positive.

\noindent\textbf{A4 (Convex calibration).}
In every active window, $K_r\ge1$, $0<\alpha_{\max}\le1$, $E_{\min}\ge1$,
$V_r\succeq0$, $\lambda_{\beta}>0$, and
$\lambda_g,\lambda_v,\lambda_s\ge0$.  The feasible set is \eqref{eq:Cr}.

\noindent\textbf{A5 (Positivity and shared support).}
Every event that enters the estimator, i.e.\ each target-supported event of
A1 that belongs to the realized active set $\mathcal{A}_r$, has strictly
positive opportunity, observation, and usable-update margins, so that the
three inversions in \eqref{eq:chi_event} are well defined there.  The ideal and
the practical reference share $\mathcal{A}_r$, $M_r$, and $\mathcal{C}_r$,
and the unnormalized H\'ajek mass obeys
$\mathbf{1}^{\top}b\ge B_{\min}>0$.  The shared-support clause of A5 fails
if Design changes $\widehat\zeta$ on $\{\pi^{\mathrm{tar}}=0\}$, which is why
RAVEN-NoDesign uses $\widehat\zeta=1$.

\noindent\textbf{A6a (Inactive stabilization at the true margins).}
On A5's active support the true margins clear every stabilizer of
\eqref{eq:design_ratio_simple}, \eqref{eq:a_stabilized}, and \eqref{eq:d_use}:
$\pi^{\mathrm{opp}}\ge\pi_{\min}$, $p^{\mathrm{obs}}\ge p_{\min}$,
$\pi^{\mathrm{tar}}/(\pi^{\mathrm{opp}}p^{\mathrm{obs}})\le a_{\max}$ with
Design on and $1/p^{\mathrm{obs}}\le a_{\max}$ with Design off, and
$1/q^{\mathrm{use}}\le d_{\max}$; when $d_{\max}\le1/q_{\min}$, as for the
frozen constants, the last clause also gives $q^{\mathrm{use}}\ge q_{\min}$.
The ceiling clause is stated on the composite weight rather than on
$1/p^{\mathrm{obs}}$: \eqref{eq:a_stabilized} already carries
$\widehat\zeta^{o}$, so bounding $1/p^{\mathrm{obs}}$ alone leaves it
unbounded.
A6a is exactly what makes the map return the ideal untruncated ratios, and
hence what makes \eqref{eq:eps_stab} vanish.

\noindent\textbf{A6b (Bounded interpolation error).}
On A5's active support the pre-outcome estimates satisfy
$|\widehat\pi^{\mathrm{opp}}-\pi^{\mathrm{opp}}|\le\varepsilon_\pi$,
$|\widehat p^{\mathrm{obs}}-p^{\mathrm{obs}}|\le\varepsilon_p$, and
$|\widehat q^{\mathrm{use}}-q^{\mathrm{use}}|\le\varepsilon_q$.

\subsection{Target Risk Is Not Arrival Risk}
\label{subsec:prop_risk}

\begin{proposition}[Target-vs-arrival risk distinction]
\label{prop:risk}
Let $\boldsymbol\varpi^{\mathrm{tar}}$ and
$\widetilde{\boldsymbol\rho}^{\mathrm{arr}}$ be two normalized laws on the same
finite supported atomic domain.  If
$\boldsymbol\varpi^{\mathrm{tar}}\neq\widetilde{\boldsymbol\rho}^{\mathrm{arr}}$,
then the two risk functionals are not interchangeable in general: there exists
a bounded nonnegative error profile $e_i$ such that
\begin{equation}
 \sum_i
 \varpi_i^{\mathrm{tar}}e_i
 \neq
 \sum_i
 \widetilde\rho_i^{\mathrm{arr}}e_i.
 \label{eq:prop1_nonidentity}
\end{equation}
Conversely, equality of the two weighted expectations for every bounded error
profile implies equality of the two underlying laws.
\end{proposition}

\noindent\emph{Proof.}
See the supplemental material.\hfill$\square$

Proposition~\ref{prop:risk} is a functional distinction, not a claim that
every predictor or minimizer must differ numerically.

\subsection{Normalized Sequential Target Reconstruction}
\label{subsec:prop_reconstruction}

For the ideal analysis, let $\mathcal{O}_r$ denote the multiset of realized
opportunity events in window $r$.  Each event $e\in\mathcal{O}_r$ carries a
mark $(k_e,i_e)$ with marginal opportunity law
$P_r^{\mathrm{opp}}(k_e,i_e)$ in \eqref{eq:opp_event_law}.  Observation and
usable-update selection then occur sequentially according to
\eqref{eq:pobs_def}--\eqref{eq:quse_def}.  Define
\begin{equation}
 \chi_e
 =
 O_e U_{k_e,r}
 \frac{
 \pi^{\mathrm{tar}}_{k_e,i_e}
 }{
 P_r^{\mathrm{opp}}(k_e,i_e)
 p^{\mathrm{obs}}_{k_e,r,i_e}
 q^{\mathrm{use}}_{k_e,r}
 }.
 \label{eq:chi_event}
\end{equation}

\begin{proposition}[Normalized target reconstruction]
\label{prop:reconstruction}
Under A1--A3, for any bounded test function $\psi(k,i)$ distinct from the
Design factor of \eqref{eq:design_factor},
\begin{equation}
 \widehat\Phi_R
 =
 \frac{
 \sum_{r<R}
 \sum_{e\in\mathcal{O}_r}
 \chi_e\psi(k_e,i_e)
 }{
 \sum_{r<R}
 \sum_{e\in\mathcal{O}_r}
 \chi_e
 }
 \xrightarrow{p}
 \Phi^{\mathrm{tar}},
 \label{eq:hajek_phi}
\end{equation}
where
\begin{equation}
 \Phi^{\mathrm{tar}}
 =
 \sum_{k,i}
 \pi^{\mathrm{tar}}_{k,i}\psi(k,i).
 \label{eq:target_phi}
\end{equation}
Therefore, client--stratum indicators recover
$\pi^{\mathrm{tar}}_{k,s}$.  By the compatibility identity
\eqref{eq:target_compatibility}, choosing
$\psi(k,i)=\mathbf{1}\{h(i)=g\}$ then recovers the deployment-group mass
$\mu_g$.  The result is moment and group-mass recovery under A1--A3, not
optimization of $\mathcal{R}_{\mu}(\theta)$.
\end{proposition}

\noindent\emph{Proof.}
See the supplemental material.\hfill$\square$

Proposition~\ref{prop:reconstruction} identifies the ideal pooled correction
with the target-mass opportunity ratio (RAVEN-Design),
not every finite-window $\boldsymbol\beta_r^\star$; that gap is why
Section~\ref{sec:calibration} calibrates inside the current support.
It does not prove that the sequential estimator optimizes deployment-target
risk.

\subsection{Existence and Uniqueness of Feasible Calibration}
\label{subsec:theorem_calibration}

\begin{theorem}[Unique solution of the calibration program]
\label{thm:calibration}
Under A4, every active window has a nonempty compact convex feasible set
$\mathcal{C}_r$, and the calibration program
$\min_{\boldsymbol\alpha\in\mathcal{C}_r}J_r(\boldsymbol\alpha)$
has a unique global minimizer $\boldsymbol\alpha_r^\star$.
\end{theorem}

\noindent\emph{Proof.}
See the supplemental material.\hfill$\square$

The theorem concerns the program, not post-hoc exact optimality of every
logged solve.

\subsection{Practical Plug-in Perturbation}
\label{subsec:prop_plugin}

Let $\mathcal{T}$ be the implemented stabilization map: the composition of
\eqref{eq:design_ratio_simple}, \eqref{eq:a_stabilized}, and \eqref{eq:d_use}
with the normalization \eqref{eq:beta_ref}.  Write
$\boldsymbol\beta_r^{\mathcal T}$ for the reference $\mathcal{T}$ returns at
the \emph{true} margins on $\mathcal{A}_r$, and $\boldsymbol\beta_r^\star$
for the ideal untruncated reference of the supplemental material.  The
stabilization error is
\begin{equation}
\varepsilon_{\mathrm{stab},r}
=
\bigl\|\boldsymbol\beta_r^{\mathcal T}-\boldsymbol\beta_r^\star\bigr\|_1 .
\label{eq:eps_stab}
\end{equation}
It vanishes exactly when $\mathcal{T}$ reproduces the ideal ratios on the
active support, and it depends only on the true margins, $\pi^{\mathrm{tar}}$
and the frozen constants, so it is measurable on a released trace.  Let
$n_r=\max_{k\in\mathcal{A}_r}\sum_i R_{k,r,i}O_{k,r,i}$ and
$K_r=|\mathcal{A}_r|$.

\begin{proposition}[Practical plug-in perturbation]
\label{prop:practical_plugin}
Assume A1, A4, A5, and A6b.  Let the stabilized window mass obey
$\mathbf{1}^{\top}b_r^{\mathcal T}\ge B_{\min}>0$ at the true margins, and let
the ideal and the implemented reference share $\mathcal{A}_r$, $M_r$, and
$\mathcal{C}_r$.  Let $\boldsymbol\beta_r^\star$ be the
same-window ideal untruncated reference and
$\widehat{\boldsymbol\beta}_r$ the implemented clipped plug-in.  Let
$\boldsymbol\alpha^\star(\boldsymbol\beta)$ denote the unique P2 minimizer
at that fixed window state.  Then
\begin{align}
\bigl\|\widehat{\boldsymbol\beta}_r-\boldsymbol\beta_r^\star\bigr\|_1
&\le
C_1\bigl(\varepsilon_\pi+\varepsilon_p+\varepsilon_q\bigr)
+\varepsilon_{\mathrm{stab},r},
\label{eq:g4_beta}\\
\bigl\|\widehat{\boldsymbol\alpha}_r-\boldsymbol\alpha_r^\star\bigr\|_2
&\le
\bigl\|\widehat{\boldsymbol\beta}_r-\boldsymbol\beta_r^\star\bigr\|_2,
\label{eq:g4_alpha}\\
\bigl\|M_r(\widehat{\boldsymbol\alpha}_r-\boldsymbol\alpha_r^\star)\bigr\|_2
&\le
\|M_r\|_2\,
\bigl\|\widehat{\boldsymbol\beta}_r-\boldsymbol\beta_r^\star\bigr\|_2,
\label{eq:g4_comp}
\end{align}
where $C_1$ is the explicit Lipschitz envelope of the stabilized H\'ajek map
(supplemental material), with
$L_\pi=\pi^{\mathrm{tar}}_{\max}/(\pi_{\min}^2p_{\min})$,
$L_p=\pi^{\mathrm{tar}}_{\max}/(\pi_{\min}p_{\min}^2)$,
$L_q=\min\{q_{\min}^{-2},d_{\max}^2\}$, and
$C_1=(2K_rn_r/B_{\min})(d_{\max}\max\{L_\pi,L_p\}+a_{\max}L_q)$; under Design
off $\zeta\equiv1$, so $L_\pi=0$ and $L_p=1/p_{\min}^2$.  The result is a
stability bound: it does not assert unbiasedness, identification on
$\{\pi^{\mathrm{tar}}=0\}$, or robustness to a change of $\mathcal{A}_r$.  It
holds with or without A6a, which is exactly what drives
$\varepsilon_{\mathrm{stab},r}$ to zero.
\end{proposition}

\noindent\emph{Proof.}
See the supplemental material.\hfill$\square$

Two corollaries fix how the bound is read.  (a) If A6a holds,
$\boldsymbol\beta_r^{\mathcal T}=\boldsymbol\beta_r^\star$ and
$\varepsilon_{\mathrm{stab},r}=0$, so \eqref{eq:g4_beta} returns to the sharp
form $C_1(\varepsilon_\pi+\varepsilon_p+\varepsilon_q)$.  (b) Under Design off
the observation side contributes nothing to $\varepsilon_{\mathrm{stab},r}$
exactly when $p^{\mathrm{obs}}\ge p_{\min}$ and
$1/p^{\mathrm{obs}}\le a_{\max}$, leaving the usable-update stage as the only
source.  The proposition is therefore not conditional on A6a, and since A6a
fails on the blocks of the main comparison this bound does not explain the
U-Air full-Design reversal.  $C_1$ is a worst-case envelope and not a numerical
certificate: at $\pi_{\min}=10^{-6}$ and $p_{\min}=0.05$ the Design-on
$L_\pi\le2\times10^{13}$, so the informative part of \eqref{eq:g4_beta} is
$\varepsilon_{\mathrm{stab},r}$.

\subsection{Aggregate Coverage Guarantees}
\label{subsec:agg_coverage_theory}

The reconstruction results above do not by themselves yield a finite-sample
coverage certificate.  Deterministic pre-window pairwise LCBs are vacuous
under hidden zero-opportunity events, and the frozen stress levels
$a\in\{0.80,0.60,0.40,0.20\}$ fall in that region
(Section~\ref{subsec:pairwise_barrier}).  In a certified regime the
false-safe probability of an empirical diagnostic Gate is bounded by
$\alpha$, by the false-safe Gate corollary in the supplemental material; and
no method can simultaneously provide nontrivial per-window certification and
universal protection against arbitrary unannounced regime shifts, by the
arbitrary-drift applicability boundary proposition in the supplemental
material.

\begin{assumption}[Local score exchangeability within a certified stable regime]
\label{ass:score-exchangeability}
Conditioned on the frozen public regime descriptor $g$, the current score
$S_t$ and the $n$ historical calibration scores are exchangeable.  Operational
regime certification does not prove this assumption; it only defines the
pre-registered subset of windows for which the assumption is posited.
\end{assumption}

\begin{theorem}[Stable-regime aggregate lower coverage]
\label{thm:agg-lcb}
Assume window $t$ is in operational state CERTIFIED, $n\ge n_{\min}=19$,
Assumption~\ref{ass:score-exchangeability} holds, $\widehat C_t=f_b(F_{t-})$
uses the frozen predictor snapshot of the current block, and the calibration
scores are strictly historical same-block scores.  Let
$q_t=S_{(k_n)}^{\mathrm{cal}}$ with $k_n$ from \eqref{eq:agg_kn}, and define
$C^{\mathrm{LCB}}_t$ by \eqref{eq:agg_CLCB}.  Then
\begin{equation}
 P\bigl(C^{\mathrm{LCB}}_t\le C_t\mid g\bigr)\ge 1-\alpha.
 \label{eq:agg_thm1}
\end{equation}
For $\alpha=0.05$ the right-hand side is $0.95$.  The theorem remains
conditional on Assumption~\ref{ass:score-exchangeability} even when the
algorithmic state is CERTIFIED.
\end{theorem}

\noindent\emph{Proof.}
See the supplemental material.\hfill$\square$

If Assumption~\ref{ass:score-exchangeability} holds, Theorem~\ref{thm:agg-lcb}
bounds the Gate in that certified regime.  Adaptive RESET/CERTIFIED switching
is not proven valid, and Algorithm~\ref{alg:window_close} does not open Design
from $\mathrm{Gate}_t$.

\section{Evaluation}
\label{sec:evaluation}

We evaluate RAVEN-MCS on risk interchangeability, RAVEN-Design versus
the RAVEN-NoDesign setting, and component roles.

\begin{itemize}
    \item \textbf{Risk interchangeability:} Does two-stage non-random
    selection create a deployment-target mismatch, including under real
    taxi opportunity?
    \item \textbf{RAVEN-Design and RAVEN-NoDesign:} On the protocol-$R$
    development table, does RAVEN-Design reduce structural mismatch, and
    how do the two Design settings trade off against each other and against
    prediction?
    \item \textbf{Component roles:} Which terms move structural
    diagnostics versus prediction, and when is Design useful as a
    reconstruction factor?
\end{itemize}

The balanced control is a no-harm check.  A risk-interchangeability experiment tests whether arrival-weighted and
deployment-target-weighted risks are interchangeable.  A T-Drive taxi
replay then checks the same four-block mismatch under real $R$ and ranks
the same family with controlled $O$ and $U$.  FedAU-Window and ObsUse-Window
remain strongest there; those rows are not mixed into
Table~\ref{tab:main_effectiveness}.
the City Scanner visit-occupancy ranking table of the supplemental material ranks the same family under real
visit occupancy and is likewise not mixed into
Table~\ref{tab:main_effectiveness}.
the Pilot~2 measured-$O$ ranking table of the supplemental material ranks Pilot~2 with measured
$O$ and controlled $U$, and is not mixed into that table either.
The main comparison is a development table: it reports both Design settings
on both blocks
under always-on protocol $R$ on station traces; it is not a jointly
measured mobile ranking.  The supplemental material reports an
intermittent-occupancy protocol table on the same station clients.  A
component-removal ablation drops one component at a time: on SensorScope,
Design moves $\Delta_{\mathrm{group}}$ while $\mathrm{RMSE}_{\mu}$ stays
essentially flat, so $\Delta_{\mathrm{group}}$ is a structural diagnostic.

\subsection{Experimental Setup}
\label{subsec:eval_setup}

\subsubsection{Datasets and target semantics}

Table~\ref{tab:datasets} summarizes the evaluated datasets.  SensorScope
\cite{ingelrest2010sensorscope,barrenetxea2019sensorscope}\footnote{\url{https://doi.org/10.5281/zenodo.2654726}}
and U-Air
\cite{zheng2013uair}\footnote{\url{http://urban-computing.com/data}}
provide the main comparison and ablation evidence and the coverage-stress audit.
Original SensorScope and U-Air traces are not claimed to contain mobile
clients.  For both, the deployment-group mapping $h(i)$ is the spatial unit
crossed with the four UTC time blocks (00--06, 06--12, 12--18, and 18--24),
giving $220$ groups on SensorScope and $144$ on U-Air, with label-free
masses.  Their opportunity strata $s(i)$ are finer: spatial identity
$\times$ four time blocks $\times$ weekday/weekend.
Accordingly, $h(i)$ and $s(i)$ retain the distinct semantics required by
Section~\ref{subsec:selection_process}.  SensorScope contains 55 station-level
spatial identities and at most 440 opportunity strata, whereas U-Air contains
36 spatial identities and at most 288 strata.

NSW-Traffic
\cite{tfnsw2024traffic}\footnote{\url{https://opendata.transport.nsw.gov.au/data/dataset/nsw-roads-traffic-volume-counts-api}}
has 43 stations, four coordinate regions, and at most 32 strata.  It
supplies the risk-interchangeability evidence and an independent four-region protocol ranking
(the NSW-Traffic four-region ranking table in the supplemental material), not mixed into
Table~\ref{tab:main_effectiveness} and not a confirmation of that table.  Cross-vintage Traffic outputs are
not used for method inference.

T-Drive
\cite{yuan2010tdrive}\footnote{\url{https://www.microsoft.com/en-us/research/publication/t-drive-trajectory-data-sample/}}
is the public Microsoft Research taxi-trajectory sample; use follows that
source's published terms.  Processed trajectories yield $R$ on $9994$
cells and $5837$ taxis, with Complete-aligned $O$ and $U$ controls and
the same four UTC groups $h(i)$.  The ranking uses the Design-off setting
there; Design-on is also reported.  FedAU-Window and ObsUse-Window remain strongest
on that hybrid ranking.

City Scanner New York {PM}$_{2.5}$ is an independent hybrid ranking
the City Scanner visit-occupancy ranking table in the supplemental material.  Occupied cells are the federated
clients, $R$ is visit occupancy, RAVEN-Design is used ($0$ of $9343$ train
units), and $O$/$U$ remain controlled.
A second City Scanner ranking uses New York Pilot~2 only
the City Scanner Pilot~2 measured-$O$ ranking table in the supplemental material.
Occupied cells remain the federated clients.
$R$ is visit occupancy on the union of kept-{PM} and {NO}$_2$ GPS streams.
$O$ is measured from finite {PM}$_{2.5}$ in $[0,1000]$
($1931$ of $14960$ occupied cell-slots unobserved).
$U$ remains a Complete-aligned control ($S_{\max}=5$).
RAVEN-Design is used ($0$ of $10141$ train occupancy units).

\begin{table*}[t]
\centering
\caption{Datasets and roles.  Targets use public or training-side
information, never test labels.  T-Drive is a hybrid ranking under real
$R$ with controlled $O$ and $U$.  NSW four-region
(the NSW-Traffic four-region ranking table in the supplemental material), City Scanner {PM}$_{2.5}$
the City Scanner visit-occupancy ranking table in the supplemental material, and City Scanner Pilot~2 measured
$O$ the City Scanner Pilot~2 measured-$O$ ranking table in the supplemental material are not mixed into
Table~\ref{tab:main_effectiveness}.}
\label{tab:datasets}
\scriptsize
\setlength{\tabcolsep}{3.2pt}
\begin{tabular*}{\textwidth}{@{\extracolsep{\fill}}lcccccc}
\toprule
Dataset & Spatial IDs & Train/Val/Test & WG blocks &
$s(i)$ & Max $s$ & Role\\
\midrule
SensorScope~\cite{ingelrest2010sensorscope} & 55 & 10285/3465/3410 & 4 &
station $\times$ block $\times$ day & 440 & main comparison; ablation\\
U-Air~\cite{zheng2013uair} & 36 & 5688/1908/1908 & 4 &
spatial ID $\times$ block $\times$ day & 288 & main comparison; ablation\\
NSW-Traffic~\cite{tfnsw2024traffic} & 43 & 18576/6192/6192 & 4 &
region $\times$ block $\times$ day & 32 & risk interchangeability; four-region protocol\\
T-Drive~\cite{yuan2010tdrive} & 9994 cells & 285659/--/74569 & 4 &
cell $\times$ block $\times$ day & -- & hybrid (real $R$)\\
City Scanner~\cite{wang2023cityscanner} & 446 cells & 9343/3290/2090 & 4 &
region $\times$ block $\times$ day & 32 & hybrid (real $R$)\\
City Scanner P2 & 416 cells & 10141/2835/1984 & 4 &
region $\times$ block $\times$ day & 32 & hybrid (meas.\ $O$)\\
\bottomrule
\end{tabular*}
\end{table*}

\subsubsection{Selection protocol and compared methods}

The main comparison and ablation condition is \emph{Complete-aligned}: protocol $R$, $O$,
and $U$ are all active and non-random in the same direction.  On
SensorScope and U-Air, protocol $R$ is always-on.  A supplemental ranking
thins that $R$ to the published T-Drive active fraction $0.425$.

These layers are not measured in every block.  Opportunity is real on T-Drive,
from taxi trajectories, and on the two City Scanner blocks, from visit
occupancy.  Observation is measured on City Scanner Pilot~2 only, and
usable-update delivery is measured in no block.  The SensorScope, U-Air, and
NSW-Traffic protocol blocks are semi-synthetic.  Their strata and target maps
come from the frozen public data, while all three selection layers are
constructed, and the tables label these layers as controls.  In those blocks
opportunity is always-on and its stratum is drawn from a joint opportunity law
with equal client mass and a factor $\exp(-1.25)$ against UTC blocks 2 and 3.
Observation is Bernoulli at rate $0.35$, reduced to $0.1925$ on the same two
blocks.  The usable-update decision is a device-tier compute and network
process that disadvantages those blocks, with staleness capped at
$S_{\max}=5$, and each client-window draws at most four risk units.  Both
construction parameters are constant within a stratum rather than varying per
unit: the observation rate depends only on the UTC block, which the
opportunity stratum already fixes, and the usable-update rate is drawn once
per client-window and shared by every risk unit in it.  The reported
predictive gains describe this construction, not the range attainable when
selection varies at the unit level.  The
compounding direction is therefore part of the construction rather than a
measured property of these traces, and the structural-alignment results of
Table~\ref{tab:main_effectiveness} are conditional on it.

The main blocks assume within-stratum exchangeability,
$\nu^{\mathrm{opp}}_{k,r}(i\mid s)=\nu^{\mathrm{tar}}_{i\mid s}$, which is the
condition under which the design ratio reduces to
\eqref{eq:design_ratio_simple}.  Paired methods share the same recorded selection trace.
The five principal same-backbone baselines are
FedAvg-Window~\cite{mcmahan2017communication},
FedAsync-Window~\cite{xie2019asynchronous},
TimeAlign-Agg (a windowed time-aware aggregator in the spirit of
\cite{wang2026federated}, not an official reproduction),
Local-H\'ajek (within-client inverse-probability / H\'ajek
normalization~\cite{horvitz1952generalization,hajek1971comment}), and
TwoStage-H\'ajek (sequential H\'ajek on opportunity and
observation).  The main table additionally
reports two windowed participation baselines, FedAU-Window and
ObsUse-Window.  FedAU-Window adapts participation debiasing
\cite{wang2024lightweight} to window-close aggregation
(the windowed-adaptation specifications in the supplemental material).
ObsUse-Window applies the observation and usable-update
corrections without Design, P2, or debt.  None is an official
reproduction.  FedCure~\cite{chen2026fedcure} is not included.  Inst-Cal and Debt-Cal
expose the transition from reconstruction to calibration and are
reported in the appendix.  Central-All, Central-Delivered, and
RAVEN-SimOracle are reference-only methods.

The main comparison and the ablation report both Design settings on both
blocks.  Algorithm~\ref{alg:window_close} takes $D_t$ as a configuration;
neither setting is derived from the data.  On SensorScope, Design is worth
$1.237\%$ of $\Delta_{\mathrm{group}}$ and $0.597\%$ of
$\Delta_{c\text{-}s}$ at unchanged prediction ($+0.005\%$ of
$\mathrm{RMSE}_{\mu}$, $0.3\sigma$); on U-Air it costs $6.886\%$ of
$\mathrm{RMSE}_{\mu}$, $6.943\%$ of $\Delta_{\mathrm{group}}$, and $7.221\%$
of tail RMSE while improving $\Delta_{c\text{-}s}$ by $4.990\%$.  All four
paired comparisons use the same 20 seeds.  The choice is therefore an
endpoint-dependent deployment decision, and is reported rather than made by
the method.  A support-compatible mixture target failed a
SensorScope no-harm check and is not a table method.

All experiments except the balanced control use 20 paired seeds (the balanced control uses five).  Those seeds capture
algorithmic and EventTrace randomness, not independent cities.  No test
result tunes the selection process, parameters, groups, or strata.

\subsubsection{Metrics and statistical policy}

Prediction is measured on the supported test units by
\[
 \mathrm{RMSE}_{\mu}
 =
 \Bigl[\sum_{i\in\mathcal{I}_{\mathrm{test}}}
 \frac{\mu^{\mathrm{eval}}_{h(i)}}{n_{h(i)}}
 \bigl(\widehat Y_i-Y_i\bigr)^2\Bigr]^{1/2},
\]
where $n_g$ is the test-unit count in group $g$, $\boldsymbol\mu$ is the
deployment-target marginal $\mu_g=\sum_{i:h(i)=g}\varpi_i^{\mathrm{tar}}$,
$\boldsymbol\mu^{\mathrm{eval}}$ is the evaluation group mass normalized on
the supported test units, and $\boldsymbol\mu^{\mathrm{blk}}$ aggregates
$\boldsymbol\mu$ to the four UTC blocks.  On SensorScope and U-Air
$\boldsymbol\mu^{\mathrm{eval}}$ is uniform on the supported test units, so
$\mathrm{RMSE}_{\mu}$ is the unweighted pooled RMSE.
The U-Air and NSW-Traffic targets are standardized by their training-split
scales ($\sigma=17.363$ and $575.776$), so those $\mathrm{RMSE}_{\mu}$ entries
are dimensionless, while SensorScope is reported in physical units.  Columns
that are ratios of two runs, such as $\Delta_{\mathrm{group}}$ and
$\Delta_{c\text{-}s}$, are invariant to this standardization.
$\mathrm{WorstGroupRMSE}$ is the largest of the four UTC-block scores, each
restricting the same weights to its own test units and renormalizing them
inside the block.  The four scores are in the per-group table of the
supplemental material.  The under-served-block score uses the block whose
$\mu^{\mathrm{blk}}_b$ most exceeds opportunity mass
(Section~\ref{subsec:eval_underserved}).
$\mathrm{Tail\_RMSE}$ applies that block functional to the frozen late-day
blocks $2$ and $3$ ($12$--$24$~UTC), fixed before the runs, not ranked by
error or arrival ratio.
For the risk-interchangeability test only, $\mathrm{RMSE}_{\rho}$ replaces the
evaluation weights by the atomic arrival weights of
\eqref{eq:arrival_distribution}.  Structural
alignment is $\Delta_{\mathrm{group}}(R)=\left\|\overline{\boldsymbol\omega}_R-\boldsymbol\mu\right\|_1$
and $\Delta_{c\text{-}s}$ in \eqref{eq:delta_cs}.  Smaller is better.

All formal method comparisons are paired by seed.  We use the paired
Wilcoxon protocol~\cite{wilcoxon1945individual} and apply Holm
correction~\cite{holm1979simple} within the predeclared comparison
families at $\alpha=0.05$; absolute differences below
$10^{-12}$ are treated as numerical ties.  Tables report
mean $\pm$ standard deviation.  On four-block $\mathrm{WorstGroupRMSE}$,
RAVEN-Design (SensorScope) and RAVEN-NoDesign (U-Air) are each compared with FedAU-Window and ObsUse-Window by a
pre-specified $+0.5\%$ mean-relative point-estimate tolerance: the mean
of the paired relative differences must be at most $+0.5\%$.
That ranking prose uses TEST seeds $30011$--$30020$; Table~\ref{tab:main_effectiveness} reports 20-seed means.
Holm tests versus those two comparators use the four 20-seed
$\mathrm{WorstGroupRMSE}$ comparisons.
This descriptive tolerance rule is not a formal non-inferiority test
and carries no coverage guarantee.  The U-Air comparison is a fallback
evaluation there, not a target-mass claim.  A later under-served-block check uses the same
$0.5\%$ mean-relative rule on validation only; a win would have required
the table RAVEN configuration to beat both FedAU-Window and ObsUse-Window by at least $0.5\%$.
For P2, warning-affected allocations satisfy
the recorded feasibility constraints; the stored outputs do not
support retrospective certification of exact numerical optimality.

\subsection{Risk Interchangeability under Two-Stage Selection}
\label{subsec:eval_rq1}

Under the balanced control, TimeAlign obtains mean deployment-target
RMSE $7.7108\pm0.0272$ and RAVEN $7.7154\pm0.0337$ (mean relative
degradation $0.0599\%$; one-sided 95\% upper bound $0.134\%$, below the
3\% no-harm margin).  The balanced control is a sanity check, not a superiority test.

We then test the premise of Proposition~1 using the NSW-Traffic risk-interchangeability
experiment.  For each method and seed, write
$\Delta_{\mu}=\mathrm{RMSE}_{\mu}^{\mathrm{aligned}}
-\mathrm{RMSE}_{\mu}^{\mathrm{balanced}}$ and
$\Delta_{\rho}=\mathrm{RMSE}_{\rho}^{\mathrm{aligned}}
-\mathrm{RMSE}_{\rho}^{\mathrm{balanced}}$, and the pre-specified
difference-of-differences $D_{\mathrm{aligned}}=\Delta_{\mu}-\Delta_{\rho}$.
If the two risk views reacted interchangeably to the same selection shift,
$D_{\mathrm{aligned}}$ would concentrate around zero.

Instead, Fig.~\ref{fig:risk_noninterchangeability} shows a stable negative
shift for all five methods.  Mean $D_{\mathrm{aligned}}$ ranges from
$-0.0209$ to $-0.0214$; 17--18 of the 20 paired seeds are negative for each
method, and the paired Wilcoxon tests against zero give
$p=4.83\times10^{-4}$--$5.86\times10^{-4}$.  The mean target-risk change is
only about $8\times10^{-4}$, whereas the mean arrival-risk change is about
$2.2\times10^{-2}$.  Thus the same two-stage selection shift changes the
arrival-weighted and deployment-target-weighted views by materially different
amounts.  Complete method-level risk-interchangeability statistics are in the risk-semantic appendix of the supplemental material.

\begin{figure}[t]
\raggedright
\includegraphics[width=\columnwidth]{fig2_e2_risk_noninterchangeability.pdf}
\caption{Risk interchangeability on complete-aligned NSW-Traffic.  (a)~Seed-level
$D_{\mathrm{aligned}}=\Delta\mathrm{RMSE}_{\mu}-\Delta\mathrm{RMSE}_{\rho}$
over 20 paired seeds.  (b)~Count of negative seeds (not a method ranking).
A negative value means that the arrival-weighted risk changes more than the
deployment-target risk under the same selection shift.  Local-H\'ajek is omitted
here because the risk-semantic experiment reports the five methods used in that test.}
\label{fig:risk_noninterchangeability}
\end{figure}

\noindent\textbf{Risk-interchangeability finding.}
The risk-interchangeability experiment supports Proposition~1: under two-stage non-random selection, the
arrival law need not be an interchangeable proxy for the deployment
target.  It is not a method-ranking test.

The same four-block mismatch appears under real taxi opportunity.
Fig.~\ref{fig:g3_tdrive} compares the public H1--H4 masses with the
mass induced by processed T-Drive trajectories.  Real $R$ is intermittent (mean client active fraction $0.425$,
burstiness $+0.38$, cell coverage $0.547$).  Four-block
$\mathrm{TV}(\pi^{\mathrm{opp}},\boldsymbol\mu)$ is $0.142$ on T-Drive,
$0.192$ on protocol SensorScope, and $0.278$ on protocol U-Air, all above
$0.05$.  $O$ and $U$ remain Complete-aligned controls.
Design is off on this block.

\begin{figure}[t]
\raggedright
\includegraphics[width=\columnwidth]{fig_g3_tdrive_h_mass.pdf}
\caption{T-Drive real-mobility opportunity.  (a)~Public $\boldsymbol\mu$,
taxi opportunity, and a unit-balanced reference; real $R$ under-serves
H3 and over-serves H4.  (b)~Four-block
$\mathrm{TV}(\pi^{\mathrm{opp}},\boldsymbol\mu)$.
Design is off; the RAVEN-NoDesign ranking is in
Table~\ref{tab:tdrive_hybrid}.}
\label{fig:g3_tdrive}
\end{figure}

\noindent\textbf{Hybrid ranking under real taxi opportunity.}
Table~\ref{tab:tdrive_hybrid} ranks the same windowed family under
real taxi $R$ with Complete-aligned $O/U$ controls.
FedAU-Window and ObsUse-Window remain strongest.
RAVEN-NoDesign does not meet the $+0.5\%$ mean-relative
tolerance versus those two (paired relative $+2.39\%$ and $+2.43\%$;
versus FedAvg $+0.32\%$).  Design on raises
$\mathrm{RMSE}_{\mu}$, worst-group RMSE, and $\Delta_{\mathrm{group}}$,
as on U-Air.  RAVEN-NoDesign leaves the four-block
structural gap unchanged relative to FedAvg.

\begin{table}[!htbp]
\centering
\caption{T-Drive hybrid ranking under real taxi $R$
(mean $\pm$ std, 10 seeds, $178$ windows).  Lower is better.
$O$ and $U$ are Complete-aligned controls.
RAVEN-NoDesign ($\widehat\zeta=1$) is not target-mass reconstruction;
RAVEN-Design (Design on) is also reported on U-Air.
WorstGroupRMSE is the maximum over UTC H1--H4.
RAVEN-NoDesign is outside the descriptive $+0.5\%$ worst-group tolerance versus
FedAU-Window and ObsUse-Window; this table is not mixed into
Table~\ref{tab:main_effectiveness}.}
\label{tab:tdrive_hybrid}
\scriptsize
\setlength{\tabcolsep}{1pt}
\begin{tabular}{lcccc}
\toprule
Method & $\mathrm{RMSE}_{\mu}$ & WorstGroupRMSE &
$\Delta_{\mathrm{group}}$ & $\Delta_{c\text{-}s}$\\
\midrule
FedAvg-Window     & \shortstack{$1.1042$\\$\pm0.0022$} & \shortstack{$1.3064$\\$\pm0.0035$} & $0.5232$ & \shortstack{$1.9758$\\$\pm0.0004$}\\
TwoStage-H\'ajek    & \shortstack{$1.1254$\\$\pm0.0142$} & \shortstack{$1.3373$\\$\pm0.0198$} & \shortstack{$0.5948$\\$\pm0.0124$} & \shortstack{$1.9642$\\$\pm0.0007$}\\
FedAU-Window      & \shortstack{$1.0898$\\$\pm0.0050$} & \shortstack{$1.2801$\\$\pm0.0109$} & $0.5232$ & \shortstack{$1.9915$\\$\pm0.0012$}\\
ObsUse-Window     & \shortstack{$1.0896$\\$\pm0.0049$} & \shortstack{$1.2796$\\$\pm0.0108$} & $0.5232$ & \shortstack{$1.9916$\\$\pm0.0012$}\\
RAVEN-NoDesign & \shortstack{$1.1069$\\$\pm0.0028$} & \shortstack{$1.3106$\\$\pm0.0044$} & $0.5232$ & \shortstack{$1.9936$\\$\pm0.0003$}\\
RAVEN-Design & \shortstack{$1.1242$\\$\pm0.0159$} & \shortstack{$1.3355$\\$\pm0.0225$} & \shortstack{$0.5948$\\$\pm0.0124$} & \shortstack{$1.9774$\\$\pm0.0006$}\\
\bottomrule
\end{tabular}
\end{table}

\noindent\textbf{Independent four-region ranking on NSW-Traffic.}
The heuristic is applied to a locally frozen four-region map
with Design on, the ranking uses RAVEN-Design.
That map is not the SensorScope/U-Air client--stratum lookup.
the NSW-Traffic four-region ranking table in the supplemental material is a protocol ranking, not mixed
into Table~\ref{tab:main_effectiveness} and not a confirmation of that
table or of a transferred support object.
TwoStage-H\'ajek is lowest on $\mathrm{RMSE}_{\mu}$ and WorstGroupRMSE.
FedAU-Window is lowest on $\Delta_{\mathrm{group}}$ ($0.0606$ versus
$0.1048$ for RAVEN-Design).  RAVEN-Design
stays within the descriptive $+0.5\%$ worst-group tolerance versus
FedAU-Window and ObsUse-Window.

\noindent\textbf{Hybrid ranking under City Scanner visit occupancy.}
the City Scanner visit-occupancy ranking table of the supplemental material ranks the same windowed family on
occupied New York {PM}$_{2.5}$ cells, with real visit occupancy $R$,
Complete-aligned $O/U$ controls, and RAVEN-Design with Design on.
The target map uses test-group shares $\mu_{h(i)}$ on training atoms, with
all four regions positive, not the station client--stratum lookup.
Paired mean relative error versus FedAU-Window and ObsUse-Window is $-1.71\%$
and $-1.34\%$ on $\mathrm{RMSE}_{\mu}$, and $-1.56\%$ and $-1.23\%$ on
worst-group RMSE.
FedAvg-Window remains lowest on both prediction endpoints
($0.2379$ and $0.2489$), while RAVEN-Design is smallest on
$\Delta_{\mathrm{group}}$ ($0.0706$).
Clients are occupied cells, not vehicles.
The ranking is not mixed into Table~\ref{tab:main_effectiveness} or
Table~\ref{tab:tdrive_hybrid}.
The target-coverage and City Scanner section of the supplemental material lists the frozen coverage statistic per block.

\noindent\textbf{Hybrid ranking under measured City Scanner observation.}
the Pilot~2 measured-$O$ ranking table of the supplemental material ranks the same windowed family on
New York Pilot~2 occupancy with measured $O$ and controlled $U$.
TwoStage-H\'ajek is lowest on $\mathrm{RMSE}_{\mu}$ and WorstGroupRMSE.
RAVEN-Design is smallest on $\Delta_{\mathrm{group}}$ and
lower than FedAU-Window, ObsUse-Window, and FedAvg-Window on both prediction endpoints
(paired mean relative $-0.43\%$, $-0.45\%$, and $-0.65\%$ on
$\mathrm{RMSE}_{\mu}$).
Clients are occupied cells, not vehicles.
The ranking is not mixed into Table~\ref{tab:main_effectiveness},
Table~\ref{tab:tdrive_hybrid}, or the City Scanner visit-occupancy ranking table of the supplemental material.

\subsection{Structural Alignment, Target RMSE, and Worst-Group Risk}
\label{subsec:eval_rq2}

Table~\ref{tab:main_effectiveness} is the protocol development table.  It
reports both Design settings on both blocks.
On SensorScope, Design reduces
$\Delta_{\mathrm{group}}$/$\Delta_{c\text{-}s}$ from FedAvg's
$1.0558/1.2903$ to $1.0139/1.2644$ ($3.96\%/2.01\%$); the Design-off row gives
$1.0265/1.2719$, so the two settings differ by $1.237\%$ and $0.597\%$
(paired, 20 seeds) while their $\mathrm{RMSE}_{\mu}$ is indistinguishable
($7.9380$ versus $7.9384$, $0.3\sigma$).  Four-block
worst-group RMSE is interchangeable with FedAU-Window and ObsUse-Window
(TEST $-0.01\%$ and $+0.05\%$; ALL20 $+0.02\%$ and $+0.10\%$).

On U-Air the two settings separate.  Design-off obtains
$\Delta_{\mathrm{group}}=1.0165$ and $\mathrm{RMSE}_{\mu}=1.1131$,
versus $1.0924$ and $1.1954$ with Design on.  Fig.~\ref{fig:w1_design_gate}
diagnoses the reversal: estimated Design and oracle Design both worsen the
primaries by about $7\%$ relative to No-Design, so the harm is
support-incompatible $\pi^{\mathrm{tar}}$, not plug-in failure.
Design still improves $\Delta_{c\text{-}s}$ on U-Air ($1.2334$ versus
$1.2949$, $4.990\%$, $17.1\sigma$), so the U-Air choice is an
endpoint-dependent trade and not a dominated setting.  On the
four-block worst-group endpoint, the TEST mean relative differences versus
FedAU-Window and ObsUse-Window are $+0.48\%$ and $+0.44\%$
(seeds $30011$--$30020$).
The same paired relatives over all $20$ table seeds are $+0.002\%$ and
$+0.075\%$.
Both slices stay within the descriptive $+0.5\%$ tolerance.  That U-Air comparison is a
fallback evaluation, not a target-mass claim.
Holm-adjusted tests on the $20$-seed $\mathrm{WorstGroupRMSE}$ family against
those two are not significant.
Within the Design-on family, RAVEN-Design lowers worst-group RMSE by $2.46\%$
against TwoStage-H\'ajek on the TEST seeds ($10/10$).  Both settings reconstruct
target mass, so this is a family-internal comparison rather than a baseline win.
Against the no-Design baselines, RAVEN-NoDesign ties all six on prediction and
on worst-group RMSE, with no Holm-significant difference in either direction.
The two settings separate on structure: RAVEN-NoDesign lowers
$\Delta_{\mathrm{group}}$ by $2.4\%$ to $3.7\%$ and $\Delta_{c\text{-}s}$ by
$1.4\%$ to $2.6\%$ against FedAvg, FedAsync, TimeAlign, FedAU and ObsUse, all
$20/20$, and ties Debt-Cal on both.  Design-on does not improve on those
baselines: its worst-group RMSE is $1.2630$ against FedAvg $1.1690$.  We do not
claim a best worst-group method.

The sealed same-backbone family, including FedAsync, TimeAlign,
Inst-Cal, and Debt-Cal, is in the supplemental material.  On U-Air
those RAVEN bars are Design-on, not the RAVEN-NoDesign rows of
Table~\ref{tab:main_effectiveness}.  Local-H\'ajek coincides with
FedAvg on U-Air because both restrict the active set to the Design support and
the realized within-client H\'ajek weights are uniform (supplemental material).

\begin{table}[!htbp]
\centering
\caption{Development table for the main comparison under always-on protocol $R$
(mean $\pm$ std, 20 seeds; lower is better).
This is a development table on station traces whose $O$ and $U$ layers are
protocol-constructed, not a
real-mobility ranking.
Both Design settings are reported on both blocks, so the table is a
within-dataset factor rather than a single configuration.
On SensorScope, Design improves $\Delta_{\mathrm{group}}$ by $1.237\%$ and
$\Delta_{c\text{-}s}$ by $0.597\%$ at unchanged prediction
($+0.005\%$ of $\mathrm{RMSE}_{\mu}$); on U-Air it costs $6.886\%$ of
$\mathrm{RMSE}_{\mu}$, $6.943\%$ of $\Delta_{\mathrm{group}}$, and
$7.221\%$ of tail RMSE while improving $\Delta_{c\text{-}s}$ by $4.990\%$
(paired, 20 seeds).
WorstGroupRMSE is the maximum over the four UTC groups.
FedAU-Window and ObsUse-Window are windowed adaptations.
On SensorScope both RAVEN settings, and on U-Air RAVEN-NoDesign, stay
within the descriptive $+0.5\%$ worst-group tolerance versus those two
(TEST seeds $30011$--$30020$); U-Air RAVEN-Design is outside it at $+7.9\%$ and $+8.0\%$.
We do not claim a Holm-significant RMSE win.
A supplemental protocol table uses intermittent occupancy.}
\label{tab:main_effectiveness}
\scriptsize
\setlength{\tabcolsep}{1pt}
\begin{tabular}{llcccc}
\toprule
Dataset & Method & $\mathrm{RMSE}_{\mu}$ & WorstGroupRMSE &
$\Delta_{\mathrm{group}}$ & $\Delta_{c\text{-}s}$\\
\midrule
\multirow{6}{*}{SensorScope}
& FedAvg-Window     & \shortstack{$7.9389$\\$\pm0.0459$} & \shortstack{$8.6716$\\$\pm0.0471$} & \shortstack{$1.0558$\\$\pm0.0288$} & \shortstack{$1.2903$\\$\pm0.0306$}\\
& TwoStage-H\'ajek    & \shortstack{$7.9309$\\$\pm0.0543$} & \shortstack{$8.6633$\\$\pm0.0556$} & \shortstack{$1.0353$\\$\pm0.0347$} & \shortstack{$1.2910$\\$\pm0.0324$}\\
& FedAU-Window      & \shortstack{$7.9364$\\$\pm0.0507$} & \shortstack{$8.6691$\\$\pm0.0522$} & \shortstack{$1.0719$\\$\pm0.0263$} & \shortstack{$1.3057$\\$\pm0.0313$}\\
& ObsUse-Window     & \shortstack{$7.9293$\\$\pm0.0513$} & \shortstack{$8.6618$\\$\pm0.0528$} & \shortstack{$1.0650$\\$\pm0.0272$} & \shortstack{$1.3062$\\$\pm0.0311$}\\
& RAVEN-Design & \shortstack{$7.9380$\\$\pm0.0473$} & \shortstack{$8.6707$\\$\pm0.0486$} & \shortstack{$\mathbf{1.0139}$\\$\pm0.0305$} & \shortstack{$\mathbf{1.2644}$\\$\pm0.0312$}\\
& RAVEN-NoDesign & \shortstack{$7.9384$\\$\pm0.0468$} & \shortstack{$8.6711$\\$\pm0.0481$} & \shortstack{$1.0265$\\$\pm0.0278$} & \shortstack{$1.2719$\\$\pm0.0308$}\\
\midrule
\multirow{6}{*}{U-Air}
& FedAvg-Window     & \shortstack{$1.1119$\\$\pm0.0603$} & \shortstack{$1.1690$\\$\pm0.0692$} & \shortstack{$1.0446$\\$\pm0.0397$} & \shortstack{$1.3138$\\$\pm0.0337$}\\
& TwoStage-H\'ajek    & \shortstack{$1.2202$\\$\pm0.0672$} & \shortstack{$1.2907$\\$\pm0.0753$} & \shortstack{$1.1196$\\$\pm0.0445$} & \shortstack{$1.2406$\\$\pm0.0297$}\\
& FedAU-Window      & \shortstack{$1.1128$\\$\pm0.0600$} & \shortstack{$1.1701$\\$\pm0.0689$} & \shortstack{$1.0557$\\$\pm0.0370$} & \shortstack{$1.3291$\\$\pm0.0360$}\\
& ObsUse-Window     & \shortstack{$1.1120$\\$\pm0.0600$} & \shortstack{$1.1692$\\$\pm0.0689$} & \shortstack{$1.0413$\\$\pm0.0374$} & \shortstack{$1.3235$\\$\pm0.0349$}\\
& RAVEN-Design & \shortstack{$1.1954$\\$\pm0.0635$} & \shortstack{$1.2630$\\$\pm0.0716$} & \shortstack{$1.0924$\\$\pm0.0479$} & \shortstack{$\mathbf{1.2334}$\\$\pm0.0320$}\\
& RAVEN-NoDesign & \shortstack{$1.1131$\\$\pm0.0638$} & \shortstack{$1.1703$\\$\pm0.0734$} & \shortstack{$\mathbf{1.0165}$\\$\pm0.0458$} & \shortstack{$1.2949$\\$\pm0.0293$}\\
\bottomrule
\end{tabular}
\end{table}

\begin{figure}[t]
\raggedright
\includegraphics[width=\columnwidth]{fig_w1_design_gate.pdf}
\caption{Design versus No-Design, estimated and oracle propensities
(20 seeds; positive $=$ Design worse).  On SensorScope, Design slightly
reduces $\Delta_{\mathrm{group}}$ ($-1.23\%$ estimated, $-1.03\%$
oracle) with a flat RMSE.  On U-Air, estimated and oracle Design both
worsen $\Delta_{\mathrm{group}}$ and $\mathrm{RMSE}_{\mu}$ by about
$7.5\%$.  Algorithm~\ref{alg:window_close} runs RAVEN-NoDesign on U-Air;
this figure diagnoses the reversal, not the
selection rule.  The U-Air Design-on $\Delta_{c\text{-}s}$ gain remains in
Table~\ref{tab:e4_ablation}.}
\label{fig:w1_design_gate}
\end{figure}

These findings agree with the scope of the theory.  Proposition~2 recovers
compatible target moments or target-group mass for the \emph{ideal pooled}
sequential correction under A1--A3 with the target-mass opportunity ratio;
it does not prove optimization of
$\mathcal{R}_{\mu}(\theta)$, and it does not claim that every finite-window
reference vector exactly equals $\boldsymbol{\mu}$.  Proposition~\ref{prop:practical_plugin}
applies when A5 and A6b hold and $\mathcal{A}_r$ is fixed, with the
stabilization error \eqref{eq:eps_stab} carried explicitly; A6a is the
additional condition that makes that error vanish.  The
support-hypothesis audit in the supplemental material reports A5 holding, A6b
untested, and A6a failing on both datasets, so on these traces
Proposition~\ref{prop:practical_plugin} is used in the form that carries
$\varepsilon_{\mathrm{stab}}$, whose value the same audit reports, rather than
in the form that cancels it.  The main comparison evaluates the
finite-window system after the feasibility-aware calibration in (P2).

\noindent\textbf{Development-table findings.}
On SensorScope, RAVEN-Design reduces structural mismatch relative to
FedAvg and stays within the $+0.5\%$ worst-group tolerance; that is not
an RMSE win.  On U-Air,
RAVEN-NoDesign improves primaries and worst-group RMSE relative to the
Design-on setting.  Against the no-Design baselines it is a tie on both
prediction and worst-group RMSE, and the structural advantage holds against
five of the six, with Debt-Cal tying on both endpoints.  Versus FedAU-Window and
ObsUse-Window,
TEST mean relative worst-group differences are $+0.48\%$ and $+0.44\%$;
target RMSE is interchangeable, not a Holm-significant win.
The under-served-block check did not produce a second external win.
Under intermittent occupancy (supplemental material), the $+0.5\%$
tolerance and the U-Air Design reversal keep their sign; the SensorScope
$\Delta_{c\text{-}s}$ gap versus FedAvg does not.

\subsection{A Pre-Registered Under-Served-Block Check}
\label{subsec:eval_underserved}

Four-block $\mathrm{WorstGroupRMSE}$ can coincide with the intrinsically
hardest UTC group rather than the group whose public mass most exceeds
opportunity mass.  After freezing the heuristic assignment, we scored a
second validation-only endpoint
$g^{\star}=\arg\max_b(\mu^{\mathrm{blk}}_b-\pi^{\mathrm{opp},4}_b)$
(SensorScope $g^{\star}=2$, gap $+0.114$; U-Air $g^{\star}=3$, gap
$+0.172$).  A win requires mean paired relative RMSE versus FedAU-Window
\emph{and} versus ObsUse-Window of at most $-0.5\%$.
Neither dataset meets that rule (SensorScope $+0.05\%/+0.15\%$;
U-Air $-0.44\%/-0.27\%$); both satisfy no-harm ($\le+0.5\%$).
The test seeds were not opened.  the under-served-block table of the supplemental material in
the supplemental material records the look.  The descriptive
$+0.5\%$ worst-group tolerance remains the only external endpoint statement.

\subsection{What the Correction and Calibration Components Do}
\label{subsec:eval_rq3}

The component-removal ablation drops Design, Obs, Use, Inst, or Debt under the same
Complete-aligned selection trace and paired seeds.
Table~\ref{tab:e4_ablation} reports the relative change
$\delta_{\mathrm{remove}}(M)
=(M_{\mathrm{ablated}}-M_{\mathrm{full}})/M_{\mathrm{full}}\times100\%$,
so a positive number means that removing the component worsens the metric;
the magnified component-removal figure in the supplemental material visualizes the
same one-factor removals.
The Inst ablation uses the one-factor removal under the same paired seeds, and
all significance statements below use the pre-specified Holm families.

\begin{table}[!htbp]
\centering
\caption{Component-removal ablation: relative change (\%) after removing one RAVEN component.
Positive means the ablation worsens the metric; negative means it improves the
metric.  $^\ast$ denotes Holm-adjusted $p<0.05$.}
\label{tab:e4_ablation}
\setlength{\tabcolsep}{3.5pt}
\begin{tabular}{llrrrr}
\toprule
Dataset & Removed & $\Delta_{\mathrm{group}}$ &
$\Delta_{c\text{-}s}$ & $\mathrm{RMSE}_{\mu}$ & Tail RMSE\\
\midrule
\multirow{5}{*}{SensorScope}
& Design & $+1.237^\ast$ & $+0.597^\ast$ & $+0.005$ & $+0.005$\\
& Obs    & $+0.096$       & $-0.070^\ast$ & $+0.012$ & $+0.011$\\
& Use    & $+0.026$       & $+0.026$       & $+0.012$ & $+0.012$\\
& Inst   & $-0.260$       & $+0.310^\ast$ & $-0.002$ & $-0.002$\\
& Debt   & $+0.014^\ast$ & $+0.004^\ast$ & $<0.001$ & $<0.001$\\
\midrule
\multirow{5}{*}{U-Air}
& Design & $-6.943^\ast$ & $+4.990^\ast$ & $-6.885^\ast$ & $-7.221^\ast$\\
& Obs    & $+0.123^\ast$ & $+0.022$       & $-0.010$       & $-0.010$\\
& Use    & $+0.014$       & $-0.014$       & $+0.025$       & $+0.026$\\
& Inst   & $-0.089$       & $-0.143$       & $+0.143^\ast$ & $+0.149^\ast$\\
& Debt   & $+0.004^\ast$ & $+0.002^\ast$ & $<0.001$ & $<0.001$\\
\bottomrule
\end{tabular}
\end{table}

\paragraph{Design correction.}
Removing Design worsens $\Delta_{c\text{-}s}$ by $0.597\%$ on SensorScope and
$4.990\%$ on U-Air, both Holm-significant.  On U-Air the same ablation
improves $\Delta_{\mathrm{group}}$ by $6.943\%$ and
$\mathrm{RMSE}_{\mu}$ by $6.886\%$.  The U-Air ranking uses the Design-off
setting, not because estimated propensities fail
(Fig.~\ref{fig:w1_design_gate}).  The $25.5\%$ figure is the realized always-on
EventTrace occurrence share ($13040$ of $51200$ appearances), not an
input of Algorithm~\ref{alg:window_close}, and is reported with the $\mathcal{A}_r$ overlap
$51.6\%$.  Oracle Design tracks estimated
Design.  SensorScope keeps Design on.  A 774-window check found identical
R/O/U and registration events, with differences beginning only at the
positive-weight/P2-support layer
(supplemental material).
The non-Design components are magnified in the U-Air ablation figure of the supplemental material.

\paragraph{Observation and usable-update correction.}
Observation correction has a smaller, endpoint-specific effect.  On U-Air,
removing Obs worsens $\Delta_{\mathrm{group}}$ by $0.123\%$
(Holm-significant), while its $\Delta_{c\text{-}s}$ and prediction effects are
not significant.  On SensorScope, the $+0.096\%$ group change is not
significant after Holm correction, whereas
$\Delta_{c\text{-}s}$ slightly improves by $0.070\%$ when Obs is removed.
Thus Obs is best characterized as another endpoint trade-off.  In contrast,
removing Use produces no significant change in any of the four core endpoints
on either dataset.  The ablation therefore does not establish the marginal
indispensability of usable-update correction in this Complete-aligned family;
the theoretical selection layer remains distinct, but its isolated empirical
effect is not clearly established here.

\paragraph{Instantaneous calibration and debt.}
Inst also exhibits a trade-off.  On SensorScope, removing it worsens
$\Delta_{c\text{-}s}$ by $0.310\%$ (significant), while the
$-0.260\%$ change in $\Delta_{\mathrm{group}}$ is not significant after Holm
correction.  On U-Air, the small structural
changes are not significant, but removing Inst worsens
$\mathrm{RMSE}_{\mu}$ and tail RMSE by $0.143\%$ and $0.149\%$,
respectively (both $p_{\mathrm{Holm}}=0.0194$).  Hence Inst exhibits an
endpoint trade-off, rather than a purely structural role.

Debt has the narrowest but most consistent role.  Removing it significantly
increases both structural mismatches on both datasets, but only by
$0.0141\%/0.0040\%$ on SensorScope and
$0.0043\%/0.0024\%$ on U-Air, well below the $0.5\%$ paired-relative
minimum practical effect used elsewhere; prediction changes are negligible.
This matches long-term accounting in \eqref{eq:debt_bound}, not a
standalone performance driver, and does not establish
$\|\mathbf Q_R\|_1=o(S_R)$.

\noindent\textbf{Component-removal finding.}
RAVEN is a coupled correction-and-calibration system, not five monotone
boosters.  As a reconstruction factor, Design helps SensorScope and
reverses U-Air primaries while still buying $\Delta_{c\text{-}s}$.
Removing Design on SensorScope worsens $\Delta_{\mathrm{group}}$ by
$1.237\%$ (Holm-significant) while $\mathrm{RMSE}_{\mu}$ changes by
$0.005\%$, so $\Delta_{\mathrm{group}}$ is not evidence that completion
became more accurate.
Obs has weaker endpoint-specific effects.  Usable-update correction is
kept as a sequential selection layer; the ablation does not show an isolated
effect.  Inst trades finite-window structure against prediction.
Target debt is a consistent accounting term whose isolated changes stay
well below the $0.5\%$ practical threshold.

\subsection{Recoverable-Gain Ceiling}
\label{subsec:ceiling}

The recoverable gain is bounded inside the aggregate class that the method
defines: weights in $\mathcal{C}_r$ of \eqref{eq:Cr}, selected by (P2), and
scored by the deployment-target criterion $\mathrm{RMSE}_{\mu}$.  Two measured
quantities bound it.  The first is the headroom from the uncorrected member
(FedAvg-Window) to the best member evaluated on true margins, reported as
$H=(\mathrm{RMSE}_{\mu}(\text{FedAvg})-\mathrm{RMSE}_{\mu}(\text{oracle}))/\mathrm{RMSE}_{\mu}(\text{FedAvg})$
over a stress grid that varies the opportunity tilt, the tail observation scale,
the model capacity and the target.  The GO rule was frozen before the runs: GO if
any condition reaches $H\ge0.5\%$ with a positive interval lower bound.  It did
not fire.  SensorScope is conclusive at a paired-difference minimum detectable
effect of $0.17$--$0.36$ percentage points, with all six estimates inside
$[-0.03\%,+0.08\%]$.  U-Air was re-run at $20$ seeds, $270$ runs and $45$
regenerated traces, where every estimate satisfies $|H|\le0.217\%$ and every
$95\%$ interval contains zero.  Its realized detectable effect is $0.70$--$0.80$
points, so that block excludes effects above about $0.7\%$ and does not resolve
the $0.5\%$ bar on its own.  Hardening the observation propensity to vary within
a stratum and to correlate with the label raises the oracle gap by
$15$--$21\times$, to $0.104\%$ on SensorScope and $0.251\%$ on U-Air, both inside
the same bound.  The gap is therefore a property of these endpoints rather
than of a regime that was made too easy.

The second quantity is the cost of estimating the margins themselves, measured
inside one calibrated member so that only the observation and usable-update
margins change.  Replacing those two margins with their true values moves
$\mathrm{RMSE}_{\mu}$ by $-0.008\%$ on SensorScope, oracle value minus calibrated
value, with a $95\%$ interval of $[-0.041\%,+0.025\%]$, and by $+0.040\%$ on
U-Air ($[-0.090\%,+0.171\%]$), paired over $20$ seeds.  On SensorScope every
member of the class we evaluate stays inside $0.08\%$ under the same
substitution.  On U-Air three of the four members stay inside $0.3\%$, while
TwoStage-H\'ajek sits $2.04\%$ above the oracle-margin member
($[1.33\%,2.75\%]$).  The substitution is not small before aggregation: the
estimated reference differs from the true-margin reference by a median of
$0.077$ in $\ell_1$ and by up to $1.15$ over the active windows, and it reaches
the aggregate composition at the same order for the members whose weights
normalize that reference.  Accurate margin estimates are therefore not what
holds the prediction endpoint flat; three of the four members are insensitive to
a large change in their own aggregation weights.
At the composition level the perturbation is already bounded: \eqref{eq:g4_comp}
limits the aggregate difference induced by the margin estimates to $\|M_r\|_2$
times the reference perturbation controlled by
Proposition~\ref{prop:practical_plugin}.  A metric-level bound of the same form
would require a transfer assumption between the training and the deployment
distributions, which Proposition~2 does not make, so the metric-level ceiling
above is measured, at the minimum detectable effects stated.

\subsection{Limitations}
\label{subsec:limitations}

RAVEN-Design and City Scanner
report Design on.  RAVEN-NoDesign reports Design off on U-Air and T-Drive
($\widehat\zeta^{o}=1$).  Both settings are reported on the blocks where the
trade-off is material.
On Table~\ref{tab:main_effectiveness}, SensorScope RAVEN-Design and U-Air
RAVEN-NoDesign stay within the descriptive $+0.5\%$ TEST tolerance versus
FedAU-Window and ObsUse-Window; we do not claim a Holm-significant RMSE win.
Table~\ref{tab:main_effectiveness} is a development table under
always-on protocol $R$ on station traces, not a real-mobility ranking.
Table~\ref{tab:tdrive_hybrid} is outside that tolerance:
FedAU-Window and ObsUse-Window remain strongest under real taxi $R$.
the NSW-Traffic four-region, City Scanner visit-occupancy, and
the Pilot~2 measured-$O$ ranking table of the supplemental material are not mixed into
Table~\ref{tab:main_effectiveness}.
The Design factor is an explicit configuration on each block's own
$\pi^{\mathrm{tar}}$ map; neither setting is derived from a data statistic,
and neither verifies positivity on target-positive pairs.
A1--A2 remain support and confounding boundaries, and A6a does not hold
on the released traces, so Proposition~\ref{prop:practical_plugin} is used
in its $\varepsilon_{\mathrm{stab}}$ form \eqref{eq:eps_stab} rather than in
the form that cancels the stabilization error.
The $25.5\%$ figure is EventTrace occurrence ($13040$ of $51200$), not
the target-coverage statistic of the supplemental material.
On a frozen coverage-stress oracle, pooled CERTIFIED validity is $0.959$
and U-Air $a=1.00$ is $0.750$.  Adaptive RESET/CERTIFIED switching is not
proven valid, and Design is not opened from $\mathrm{Gate}_t$.
Federated learning itself does not provide a privacy guarantee.


\section{Conclusion}
\label{sec:conclusion}

Windowed-asynchronous federated MCS completion should align selected mass
with the public deployment target, not with the arrival law of usable
updates.  RAVEN-Design reconstructs compatible target mass and
calibrates that reference inside the current feasible window.
RAVEN-NoDesign uses $\widehat\zeta=1$ and does not reconstruct target mass on that
layer.  When the two laws differ, their risks are not
interchangeable.  The ideal sequential correction recovers compatible
target moments or group mass under A1--A3; it does not prove optimization
of deployment-target risk.  Calibration is uniquely solvable, and the
clipped plug-in is a stable perturbation of the same-window reference
when target support is compatible.

On the protocol development table, SensorScope RAVEN-Design reports
structural alignment and a descriptive $+0.5\%$ worst-group tolerance,
not an RMSE win.  U-Air and T-Drive report RAVEN-NoDesign.  Hybrid tables
are mobility ceilings, not confirmations of
Table~\ref{tab:main_effectiveness}: FedAU-Window and ObsUse-Window remain
strongest on T-Drive (Table~\ref{tab:tdrive_hybrid}), FedAvg-Window is
lowest on the City Scanner visit-occupancy ranking, and TwoStage-H\'ajek is
lowest on Pilot~2 measured $O$; both rankings are reported in the
supplemental material.
$O$ and $U$ are protocol-constructed except Pilot~2 measured $O$.
Usable-update correction has no isolated effect in the tested family;
target debt remains internal accounting.
No ranking jointly measures $R$, $O$, and $U$.

\section*{Data and Code Availability}
SensorScope, U-Air, NSW-Traffic, and T-Drive are public and cited in
Section~\ref{subsec:eval_setup}.  T-Drive use follows the published
Microsoft Research terms~\cite{yuan2010tdrive}.
City Scanner New York {PM}$_{2.5}$~\cite{wang2023cityscanner}
(Zenodo 10.5281/zenodo.7961851, CC BY 4.0) is ranked in
the City Scanner visit-occupancy ranking table of the supplemental material and
the Pilot~2 measured-$O$ ranking table of the supplemental material.
The bundle \texttt{RAVEN\_MCS\_PUBLIC\_REPRO\_R1} contains paired seeds,
frozen targets, EventTraces, tabulated figure and table sources,
\texttt{SCD\_PHI.csv}, canonical Table~II
(\texttt{TABLE\_II\_MAIN\_EFFECTIVENESS.csv}),
the $20$-seed $\mathrm{WorstGroupRMSE}$ Holm family versus FedAU-Window and
ObsUse-Window, ALL20/TEST10 slice relatives, T-Drive hybrid ranking
CSVs, the NSW four-region files, City Scanner visit-occupancy ranking
CSVs (\texttt{cityscanner/}), and Pilot~2 measured-$O$ ranking CSVs
(\texttt{cityscanner\_p2/}).
A verifier reconstructs tabulated cells from those CSVs; the
bundle is not a pinned retraining archive, and NSW four-region is not a
confirmation experiment.

\bibliographystyle{IEEEtran}
\bibliography{reference}

\end{document}
